\documentclass[12pt,a4paper]{article}
\usepackage{fontspec}
\setmainfont{NotoSerif}[Path=fonts/,Extension=.ttf,UprightFont=*-Regular,BoldFont=*-Bold,ItalicFont=*-Italic,BoldItalicFont=*-BoldItalic]
\usepackage[russian]{babel}
\usepackage{amsmath,amssymb,amsthm}
\usepackage{unicode-math}
\setmathfont[Scale=1.05]{latinmodern-math.otf}
\usepackage[margin=24mm]{geometry}
\usepackage{enumitem}
\usepackage[normalem]{ulem}
\usepackage{needspace}
\usepackage[unicode,hidelinks]{hyperref}
\usepackage{microtype}
\setlength{\parindent}{0pt}
\setlength{\parskip}{6pt}
\setlength{\emergencystretch}{2em}
\linespread{1.16}\sloppy
\everymath{\displaystyle}
\theoremstyle{definition}
\newtheorem{definition}{Определение}[section]
\theoremstyle{plain}
\newtheorem{theorem}{Теорема}[section]
\theoremstyle{remark}
\newtheorem*{example}{Пример}
\newcommand{\term}[1]{\uline{#1}}
\newcommand{\starmark}{\ensuremath{{}^{\ast}}}
\title{Алгебра перед математическим анализом\\[8pt]\large Сборник задач и кратких теоретических справок}
\author{}
\date{}
\begin{document}
\maketitle
\begin{abstract}
Сборник содержит 269 задач по 14 темам: вычисления, алгебра, функции, тригонометрия и начала доказательств. В каждой главе даны определения и основные свойства, учебные задачи, задачи из жизни и три задания повышенной сложности, отмеченные звёздочкой. Ответы и краткие решения приведены в конце.
\end{abstract}
\subsection*{Как работать со сборником}
Решайте задачи по порядку, записывая промежуточные действия. В заданиях на построение рисуйте график, в заданиях на доказательство приводите рассуждение. При вычислении дробей приводите ответ к несократимому виду. В уравнениях находите все допустимые корни и проверяйте их; в неравенствах указывайте всё множество решений. Сверяйтесь с ответами после самостоятельной попытки. Звёздочка \starmark{} означает повышенную сложность относительно материала главы.
\newpage
\tableofcontents
\newpage

\clearpage
\section{Числа и свойства арифметических действий}
\subsection*{Краткая справка}
\subsubsection*{Числа и их обозначения}

\begin{definition}
\term{Натуральные числа} --- числа, используемые при счёте отдельных предметов:
\[
\mathbb N=\{1,2,3,\ldots\}.
\]
В этом сборнике отсчёт натуральных чисел начинается с \(1\).
\end{definition}

\begin{definition}
\term{Ноль} --- число \(0\), которое выражает отсутствие предметов при счёте.
Натуральные числа вместе с нулём образуют множество
\[
\mathbb N_0=\{0,1,2,3,\ldots\}.
\]
\end{definition}

\begin{definition}
\term{Единица} --- число \(1\), соответствующее одному предмету; это наименьшее натуральное число.
Ноль и единица различны: \(0\ne1\).
\end{definition}

\begin{definition}
\term{Целые числа} --- натуральные числа, противоположные им отрицательные числа и ноль:
\[
\mathbb Z=\{\ldots,-3,-2,-1,0,1,2,3,\ldots\}.
\]
\end{definition}

\begin{definition}
\term{Противоположное число} к числу \(a\) --- число \(-a\), сумма которого с \(a\) равна нулю:
\[
a+(-a)=0.
\]
Например, \(5\) и \(-5\) противоположны; число, противоположное нулю, --- сам ноль.
\end{definition}

На числовой прямой число тем больше, чем правее находится его точка.
Например, \(-8<-3<0<4\).
Запись \(a\in\mathbb Z\) читается: «\(a\) принадлежит множеству целых чисел».

\subsubsection*{Действия с целыми числами}

Сумма, разность и произведение целых чисел всегда являются целыми числами.
Деление может вывести за пределы целых чисел: например, \(6:3=2\), но \(1:2\) не является целым числом.

Вычитание означает прибавление противоположного числа:
\[
a-b=a+(-b),\qquad a-(-b)=a+b.
\]
При сложении двух отрицательных чисел складывают их расстояния до нуля и ставят минус:
\[
(-7)+(-5)=-12.
\]
При сложении чисел разных знаков из большего расстояния до нуля вычитают меньшее
и берут знак числа, расположенного дальше от нуля:
\[
-14+9=-5,\qquad -9+14=5,\qquad -9+9=0.
\]
При умножении или делении двух ненулевых чисел одинаковые знаки дают плюс, разные --- минус:
\[
(-6)(-4)=24,\qquad (-6)\cdot4=-24,\qquad
\frac{-24}{-6}=4,\qquad \frac{24}{-6}=-4.
\]
Деление на ноль не определено; в частности, запись \(\frac{0}{0}\) также не имеет значения.

\subsubsection*{От целых чисел к рациональным}

\begin{definition}
\term{Рациональное число} --- число, которое можно записать в виде
\[
\frac pq,\qquad p\in\mathbb Z,\quad q\in\mathbb Z,\quad q\ne0.
\]
Множество рациональных чисел обозначают \(\mathbb Q\).
\end{definition}

Любое целое число рационально, поскольку \(n=\frac n1\).
Например, \(3\), \(-2\), \(0\), \(\frac12\) и \(-\frac34\) --- рациональные числа.
Правила действий с дробями подробно разобраны в главе 3.
Следующие свойства удобно сформулировать сразу для всех рациональных чисел:
ими пользуются и при вычислениях, и при дальнейших преобразованиях выражений с буквами.

\subsubsection*{Основные свойства арифметических действий}

Во всех формулах ниже \(a,b,c\) --- произвольные рациональные числа.

\textbf{Замкнутость сложения и умножения.}
Сумма и произведение рациональных чисел снова рациональны:
\[
a+b\in\mathbb Q,\qquad ab\in\mathbb Q.
\]
Разность также рациональна; частное рационально при ненулевом делителе.

\textbf{Коммутативность} (переместительный закон).
Перестановка слагаемых или множителей не меняет результата:
\[
a+b=b+a,\qquad ab=ba.
\]

\textbf{Ассоциативность} (сочетательный закон).
Группировка слагаемых или множителей не меняет результата:
\[
(a+b)+c=a+(b+c),\qquad (ab)c=a(bc).
\]

\textbf{Нейтральные элементы.}
Прибавление нуля и умножение на единицу не меняют число:
\[
a+0=0+a=a,\qquad a\cdot1=1\cdot a=a.
\]
Поэтому \(0\) называют нейтральным элементом сложения,
а \(1\) --- нейтральным элементом умножения.

\textbf{Противоположное число.}
У каждого числа есть противоположное:
\[
a+(-a)=(-a)+a=0.
\]

\begin{definition}
\term{Обратное число} к ненулевому числу \(a\) --- число \(\frac1a\),
произведение которого с \(a\) равно единице:
\[
a\cdot\frac1a=\frac1a\cdot a=1,\qquad a\ne0.
\]
\end{definition}
Например, \(2\) и \(\frac12\) взаимно обратны.
У нуля обратного числа нет. Деление заменяется умножением на обратное число:
\[
\frac ab=a\cdot\frac1b,\qquad b\ne0.
\]

\textbf{Дистрибутивность} (распределительный закон).
Умножение распределяется относительно сложения:
\[
a(b+c)=ab+ac,\qquad (a+b)c=ac+bc.
\]
Для разности соответственно \(a(b-c)=ab-ac\).

\textbf{Умножение на ноль и правила знаков.}
Из перечисленных свойств следуют равенства
\[
a\cdot0=0,\qquad (-a)b=-(ab),\qquad (-a)(-b)=ab.
\]

\textbf{Законы сокращения.}
Прибавление одного числа к обеим частям равенства можно отменить;
умножение можно отменить, если общий множитель ненулевой:
\[
a+c=b+c\ \Longrightarrow\ a=b,
\]
\[
ac=bc,\quad c\ne0\ \Longrightarrow\ a=b.
\]
Знак \(\Longrightarrow\) означает «из этого следует».

\textbf{Свойство нулевого произведения.}
Произведение равно нулю тогда и только тогда, когда хотя бы один множитель равен нулю:
\[
ab=0\quad\Longleftrightarrow\quad a=0\ \text{или}\ b=0.
\]
Знак \(\Longleftrightarrow\) означает, что утверждения равносильны: каждое следует из другого.

Коммутативность и ассоциативность относятся к сложению и умножению.
Переносить их на вычитание и деление нельзя:
\[
7-2\ne2-7,\qquad (12:3):2\ne12:(3:2).
\]

\subsubsection*{Порядок действий}

\begin{enumerate}[label=\arabic*.,leftmargin=2em,itemsep=3pt]
\item Действия в скобках, начиная с самых внутренних.
\item Возведение в степень. Например, \(3^2=3\cdot3\).
\item Умножение и деление --- слева направо.
\item Сложение и вычитание --- слева направо.
\end{enumerate}

Степень относится только к своему основанию:
\[
(-3)^2=9,\qquad -3^2=-(3^2)=-9.
\]
Например,
\[
-3-(-5)\cdot2=-3-(-10)=-3+10=7.
\]

\subsubsection*{Порядок на числовой прямой}
Для любых рациональных \(a,b\) выполняется ровно одно из условий: \(a<b\), \(a=b\), \(a>b\).
Если \(a<b\) и \(b<c\), то \(a<c\) — это \textbf{транзитивность}.
Из \(a<b\) следует \(a+c<b+c\). Если также \(c>0\), то \(ac<bc\); если \(c<0\), то \(ac>bc\).
Эти правила объясняют преобразования неравенств в главе 10.
\Needspace{6\baselineskip}\subsection*{Задачи}
\begin{enumerate}[label=\thesection.\arabic*.,leftmargin=2.8em,itemsep=12pt]
\Needspace{3\baselineskip}\item Сравни: \(- 8\) и \(- 3;\allowbreak 0\) и \(- 5;\allowbreak - 4\) и \(4\).
\Needspace{3\baselineskip}\item Вычисли: \(- 7 + (- 5)\).
\Needspace{3\baselineskip}\item Вычисли: \(- 14 + 9 + 8\).
\Needspace{3\baselineskip}\item Вычисли: \(6 - (- 4) - 13\).
\Needspace{3\baselineskip}\item Вычисли: \((- 6) \cdot  (- 4) : (- 3)\).
\Needspace{3\baselineskip}\item Вычисли: \(0 \cdot  (- 15) + 0 : (- 3) + 7\). Можно ли вычислить \(7 : 0\)?
\Needspace{3\baselineskip}\item Вычисли: \(18 - 4 \cdot  7 + 12 : (- 3)\).
\Needspace{3\baselineskip}\item Вычисли: \(48 : (- 6) \cdot  2 - 9 + 5\).
\Needspace{3\baselineskip}\item Вычисли: \(- 3 \cdot  [8 - (5 - 9)] + 10\).
\Needspace{3\baselineskip}\item Вычисли: \((- 3)^{2} - 3^{2} + (- 2)^{3} - (- 2^{2})\).
\Needspace{3\baselineskip}\item Из жизни. Утром было \(- 6 {}^{\circ} C\). К полудню потеплело на \(9 {}^{\circ} C\), а вечером похолодало на \(5 {}^{\circ} C\). Какая температура стала вечером?
\Needspace{3\baselineskip}\item Из жизни. Лифт находился на уровне \(- 2\) подземной парковки, поднялся на \(7\) уровней и опустился на \(3\). На каком уровне он оказался? Уровень земли обозначен \(0\), уровни идут подряд.
\Needspace{3\baselineskip}\item Из жизни. На игровом счёте было \(120\) очков. За каждую из трёх ошибок сняли \(15\) очков, затем начислили \(40\) очков. Сколько очков осталось?
\Needspace{3\baselineskip}\item[\textbf{1.14.}\starmark] Вычисли: \(- 48 : [6 - 2 \cdot  (- 1)] + 3 \cdot  [- 4 - (- 7)]\).
\Needspace{3\baselineskip}\item[\textbf{1.15.}\starmark] Подставь скобки в \(18 - 6 : 3 + 1\), сохранив порядок чисел и действий, чтобы получилось \(3\).
\Needspace{3\baselineskip}\item[\textbf{1.16.}\starmark] В викторине за верный ответ дают \(5\) очков, за неверный снимают \(3\). Участник ответил на все \(12\) вопросов и набрал \(28\) очков. Сколько было верных ответов?
\end{enumerate}
\clearpage
\section{Делимость, простые числа, НОД и НОК}
\subsection*{Краткая справка}
\subsubsection*{Делимость и кратность}

\begin{definition}
\term{Делимость нацело}: целое число \(b\) делится на целое число \(a\ne0\),
если существует целое число \(k\), для которого
\[
b=ak.
\]
При этом \(a\) называют \term{делителем} числа \(b\).
\end{definition}

Запись с тремя вертикальными точками \(b\mathrel{\vdots}a\) читается:
«\(b\) делится на \(a\) нацело».
В университетских курсах также используют запись \(a\mid b\): «\(a\) делит \(b\)».
Порядок чисел в этих двух обозначениях противоположен:
\[
12\mathrel{\vdots}4\quad\Longleftrightarrow\quad4\mid12,
\qquad 12=4\cdot3.
\]
Здесь знак \(\Longleftrightarrow\) означает, что записи выражают одно и то же утверждение.

\begin{definition}
\term{Кратное} числа \(a\ne0\) --- целое число \(b\), которое можно представить как \(b=ak\)
с целым \(k\). Иначе говоря, \(b\) кратно \(a\), если \(b\mathrel{\vdots}a\).
\end{definition}

Например, положительные кратные числа \(4\): \(4,8,12,16,\ldots\).
Ноль также кратен любому ненулевому целому числу: \(0=a\cdot0\).
В задачах на НОК рассматриваются только \emph{положительные} общие кратные.

\subsubsection*{Простые и составные числа}

\begin{definition}
\term{Простое число} --- натуральное число, имеющее ровно два положительных делителя:
\(1\) и само это число.
\end{definition}

\begin{definition}
\term{Составное число} --- натуральное число больше \(1\), имеющее положительный делитель,
отличный от \(1\) и от самого числа.
\end{definition}

Число \(7\) простое: его положительные делители --- \(1\) и \(7\).
Число \(9\) составное: кроме \(1\) и \(9\), оно делится на \(3\).
Число \(1\) не является ни простым, ни составным.
Число \(2\) --- единственное чётное простое число.

\subsubsection*{Наибольший общий делитель и наименьшее общее кратное}

\begin{definition}
\term{Наибольший общий делитель} положительных целых чисел \(a\) и \(b\) ---
наибольшее положительное целое число, на которое оба числа делятся нацело.
Обозначения: \(\text{НОД}(a,b)\) и \(\gcd(a,b)\).
Английское название --- \emph{greatest common divisor}, сокращённо \emph{gcd}.
\end{definition}

Общие положительные делители чисел \(12\) и \(18\) --- \(1,2,3,6\), поэтому
\[
\text{НОД}(12,18)=\gcd(12,18)=6.
\]

\begin{definition}
\term{Наименьшее общее кратное} положительных целых чисел \(a\) и \(b\) ---
наименьшее положительное целое число, которое делится нацело и на \(a\), и на \(b\).
Обозначения: \(\text{НОК}(a,b)\) и \(\operatorname{lcm}(a,b)\).
Английское название --- \emph{least common multiple}, сокращённо \emph{lcm}.
\end{definition}

Положительные общие кратные чисел \(12\) и \(18\) --- \(36,72,108,\ldots\), поэтому
\[
\text{НОК}(12,18)=\operatorname{lcm}(12,18)=36.
\]

\begin{definition}
\term{Взаимно простые числа} --- положительные целые числа, НОД которых равен \(1\).
\end{definition}
Например, \(8\) и \(15\) взаимно просты, хотя оба числа составные.

\subsubsection*{Основная теорема арифметики (ОТА)}

\begin{theorem}
Каждое натуральное число больше \(1\) можно представить произведением простых чисел.
Такое разложение единственно с точностью до перестановки множителей.
\end{theorem}

Например,
\[
60=2\cdot2\cdot3\cdot5=2^2\cdot3\cdot5.
\]
Множители \(2,3,5\) простые; запись \(60=4\cdot15\) ещё не является разложением на простые множители.

\subsubsection*{Вычисление НОД и НОК по разложению}

Разложите каждое число на простые множители, объединив повторения в степени.
Для НОД возьмите только общие простые множители в \emph{меньших} степенях.
Для НОК возьмите все встречающиеся простые множители в \emph{больших} степенях.
Если общих простых множителей нет, НОД равен \(1\).

Например,
\[
12=2^2\cdot3,\qquad18=2\cdot3^2,
\]
\[
\gcd(12,18)=2\cdot3=6,\qquad
\operatorname{lcm}(12,18)=2^2\cdot3^2=36.
\]
Для нескольких положительных чисел правило такое же:
в НОД входят множители, общие \emph{всем} числам, с наименьшими показателями;
в НОК --- все встретившиеся множители с наибольшими показателями.

\begin{theorem}
Для любых положительных целых чисел \(a,b\)
\[
\gcd(a,b)\operatorname{lcm}(a,b)=ab.
\]
\end{theorem}
Это равенство относится к двум числам; напрямую переносить его на три числа нельзя.

\subsubsection*{Алгоритм Евклида}

Для нахождения НОД можно последовательно делить с остатком,
каждый раз заменяя пару «делимое, делитель» парой «делитель, остаток».
Когда остаток стал нулём, НОД равен последнему ненулевому делителю.
Например,
\[
48=2\cdot18+12,\qquad
18=1\cdot12+6,\qquad
12=2\cdot6+0,
\]
поэтому \(\gcd(48,18)=6\).

\subsubsection*{Признаки делимости}

В десятичной записи натурального числа действуют следующие признаки.
Каждый из них даёт необходимое и достаточное условие: число делится на указанный делитель
\emph{тогда и только тогда}, когда выполнена соответствующая проверка.

\begin{itemize}[leftmargin=1.5em,itemsep=4pt]
\item \textbf{На \(2\):} последняя цифра --- \(0,2,4,6\) или \(8\).
\item \textbf{На \(3\):} сумма цифр делится на \(3\).
\item \textbf{На \(4\):} число, образованное двумя последними цифрами, делится на \(4\).
\item \textbf{На \(5\):} последняя цифра --- \(0\) или \(5\).
\item \textbf{На \(6\):} число делится одновременно на \(2\) и на \(3\).
\item \textbf{На \(8\):} число, образованное тремя последними цифрами, делится на \(8\).
\item \textbf{На \(9\):} сумма цифр делится на \(9\).
\item \textbf{На \(10\):} последняя цифра --- \(0\).
\end{itemize}

Если число содержит меньше двух или трёх цифр, в признаках на \(4\) и \(8\) проверяют само число.
Например, \(123\mathrel{\vdots}3\), поскольку \(1+2+3=6\),
но \(123\) не делится на \(9\), поскольку \(6\) не делится на \(9\).
\Needspace{6\baselineskip}\subsection*{Задачи}
\begin{enumerate}[label=\thesection.\arabic*.,leftmargin=2.8em,itemsep=12pt]
\Needspace{3\baselineskip}\item Выпиши все положительные делители \(36\).
\Needspace{3\baselineskip}\item Выпиши все положительные делители \(48\).
\Needspace{3\baselineskip}\item Найди общие положительные делители \(20\) и \(30\).
\Needspace{3\baselineskip}\item Выпиши кратные \(7\) между \(20\) и \(60\).
\Needspace{3\baselineskip}\item Найди наименьшее трёхзначное число, кратное \(8\).
\Needspace{3\baselineskip}\item Раздели на простые и составные: \(17,\allowbreak 25,\allowbreak 31,\allowbreak 39,\allowbreak 49,\allowbreak 51,\allowbreak 53\).
\Needspace{3\baselineskip}\item Выпиши простые числа от \(20\) до \(40\).
\Needspace{3\baselineskip}\item Может ли простое число оканчиваться на \(0\)? Объясни.
\Needspace{3\baselineskip}\item Любое ли нечётное число простое? Приведи контрпример.
\Needspace{3\baselineskip}\item Какие числа делятся на \(3\), а какие на \(9\): \(126,\allowbreak 234,\allowbreak 405,\allowbreak 517,\allowbreak 702\)?
\Needspace{3\baselineskip}\item Какие числа делятся на \(2\) и на \(5\) одновременно: \(125,\allowbreak 230,\allowbreak 342,\allowbreak 480,\allowbreak 705\)?
\Needspace{3\baselineskip}\item Какие цифры можно поставить вместо \(\square \) в \(43 \square \), чтобы число делилось на \(3\)?
\Needspace{3\baselineskip}\item Какая цифра вместо \(\square \) в \(7 \square  2\) даёт делимость на \(9\)?
\Needspace{3\baselineskip}\item Какие цифры вместо \(\square \) в \(12 \square \) дают делимость на \(2\) и на \(3\)?
\Needspace{3\baselineskip}\item Разложи \(72\) на простые множители.
\Needspace{3\baselineskip}\item Разложи \(150\) на простые множители.
\Needspace{3\baselineskip}\item Разложи \(252\) на простые множители.
\Needspace{3\baselineskip}\item Заверши разложение \(90 = 2 \cdot  5 \cdot  9\) на простые множители.
\Needspace{3\baselineskip}\item Найди НОД и НОК чисел \(24\) и \(36\).
\Needspace{3\baselineskip}\item Найди НОД и НОК чисел \(18\) и \(45\).
\Needspace{3\baselineskip}\item Найди НОД и НОК чисел \(8\) и \(15\).
\Needspace{3\baselineskip}\item Найди НОД и НОК чисел \(12,\allowbreak 18\) и \(30\).
\Needspace{3\baselineskip}\item Из \(24\) карандашей и \(36\) ручек собери максимум одинаковых наборов, использовав всё. Сколько наборов и что в каждом?
\Needspace{3\baselineskip}\item Автобусы отправились вместе в 9:00. Интервалы отправления: \(12\) и \(18\) минут. Когда следующее совместное отправление?
\Needspace{3\baselineskip}\item Прямоугольник \(60 \times  84\) см разрежь без остатка на одинаковые квадраты наибольшей стороны. Найди сторону и число квадратов.
\Needspace{3\baselineskip}\item Выпиши все положительные делители \(24\).
\Needspace{3\baselineskip}\item Выпиши первые пять положительных кратных \(6\).
\Needspace{3\baselineskip}\item Какие из чисел \(1,\allowbreak 2,\allowbreak 9,\allowbreak 13,\allowbreak 21,\allowbreak 29\) простые?
\Needspace{3\baselineskip}\item На какие числа из списка \(2,\allowbreak 3,\allowbreak 5,\allowbreak 9,\allowbreak 10\) делится \(450\)?
\Needspace{3\baselineskip}\item Разложи \(84\) на простые множители.
\Needspace{3\baselineskip}\item Найди НОД и НОК чисел \(12\) и \(18\).
\Needspace{3\baselineskip}\item Из жизни. Для праздника есть \(42\) яблока и \(56\) мандаринов. Их нужно разложить без остатка в максимальное число одинаковых подарков. Сколько подарков получится и что будет в каждом?
\Needspace{3\baselineskip}\item Из жизни. Два индикатора загорелись одновременно. Один загорается каждые \(8\) секунд, другой каждые \(12\) секунд. Через сколько секунд они впервые снова загорятся вместе?
\Needspace{3\baselineskip}\item Из жизни. В коробке \(97\) деталей. Можно ли разложить их в одинаковые группы больше одной детали, получив больше одной группы, без остатка?
\Needspace{3\baselineskip}\item[\textbf{2.35.}\starmark] Найди наименьшее натуральное число, которое при делении на \(6,\allowbreak 8\) и \(9\) даёт остаток \(1\) и при этом больше \(1\).
\Needspace{3\baselineskip}\item[\textbf{2.36.}\starmark] НОД двух положительных целых чисел равен \(6\), НОК равен \(180\). Одно число \(30\). Найди второе.
\Needspace{3\baselineskip}\item[\textbf{2.37.}\starmark] Найди все цифры \(a\) и \(b\), чтобы четырёхзначное число \(\overline{3a5b}\) делилось на \(5\) и на \(9\).
\end{enumerate}
\clearpage
\section{Обыкновенные дроби}
\subsection*{Краткая справка}
\begin{definition}\term{Обыкновенная дробь} \(\frac{a}{b}\) — запись частного \(a : b\), где \(a\) и \(b\) — целые числа и \(b \ne  0\).\end{definition}
\begin{definition}\term{Числитель} дроби \(\frac{a}{b}\) — число \(a\), записанное над дробной чертой.\end{definition}
\begin{definition}\term{Знаменатель} дроби \(\frac{a}{b}\) — ненулевое число \(b\), записанное под дробной чертой.\end{definition}
\begin{definition}\term{Правильная положительная дробь} — дробь с положительными числителем и знаменателем, у которой числитель меньше знаменателя.\end{definition}
\begin{definition}\term{Неправильная положительная дробь} — дробь с положительными числителем и знаменателем, у которой числитель не меньше знаменателя.\end{definition}
\begin{definition}\term{Смешанное число} — сумма целой части и правильной дробной части; например, \(2 \frac{1}{3} = 2 + \frac{1}{3}\).\end{definition}
\subsubsection*{Основное свойство дроби}
\(\frac{a}{b} = \frac{ak}{bk}\) при \(k \ne  0\). Числитель и знаменатель можно умножать или делить на одно и то же ненулевое число.\par
\subsubsection*{Действия с дробями}
Сложение: \(\frac{a}{b} + \frac{c}{d} = \frac{ad + bc}{bd}\), где \(b,\allowbreak d \ne  0\).\par
Вычитание: \(\frac{a}{b} - \frac{c}{d} = \frac{ad - bc}{bd}\), где \(b,\allowbreak d \ne  0\).\par
Умножение: \((\frac{a}{b}) \cdot  (\frac{c}{d}) = \frac{ac}{bd}\).\par
Деление: \((\frac{a}{b}) : (\frac{c}{d}) = \frac{ad}{bc}\); дополнительно \(c \ne  0\).\par
Сравнение при \(b,\allowbreak d > 0\): сравни \(a \cdot  d\) и \(c \cdot  b\).\par
{\small\textit{Пример: \(\frac{1}{2} + \frac{1}{3} = \frac{3}{6} + \frac{2}{6} = \frac{5}{6}\). Сокращение: \(\frac{18}{24} = \frac{3}{4}\).}}\par
\subsubsection*{Связь со следующими темами}
Правила дробей позже применяются к выражениям с буквами; ограничения на знаменатели при этом сохраняются.
\Needspace{6\baselineskip}\subsection*{Задачи}
\begin{enumerate}[label=\thesection.\arabic*.,leftmargin=2.8em,itemsep=12pt]
\Needspace{3\baselineskip}\item Сократи \(\frac{18}{24}\).
\Needspace{3\baselineskip}\item Преврати \(\frac{17}{5}\) в смешанное число.
\Needspace{3\baselineskip}\item Преврати \(2 \frac{3}{7}\) в неправильную дробь.
\Needspace{3\baselineskip}\item Сравни \(\frac{5}{8}\) и \(\frac{7}{12}\).
\Needspace{3\baselineskip}\item Вычисли \(\frac{3}{4} - \frac{5}{6}\).
\Needspace{3\baselineskip}\item Вычисли \(\frac{7}{15} + \frac{2}{9}\).
\Needspace{3\baselineskip}\item Вычисли \((- \frac{3}{7}) \cdot  (\frac{14}{9})\).
\Needspace{3\baselineskip}\item Вычисли \((\frac{2}{3}) : (\frac{4}{9})\).
\Needspace{3\baselineskip}\item Вычисли \(2 \frac{1}{3} - 1 \frac{5}{6}\).
\Needspace{3\baselineskip}\item Вычисли \((\frac{1}{2} + \frac{1}{3}) : (\frac{5}{6})\).
\Needspace{3\baselineskip}\item Вычисли \(\frac{\frac{3}{4}}{\frac{9}{8}}\).
\Needspace{3\baselineskip}\item Из \(80\) страниц прочитали \(\frac{3}{5}\). Сколько страниц осталось?
\Needspace{3\baselineskip}\item Из жизни. На пирог ушло \(\frac{3}{4}\) кг муки, а на блины \(\frac{2}{3}\) кг. Сколько всего муки израсходовали?
\Needspace{3\baselineskip}\item Из жизни. В бутылке \(\frac{3}{2}\) л сока. Его разливают по стаканам вместимостью \(\frac{1}{4}\) л. Сколько полных стаканов получится?
\Needspace{3\baselineskip}\item Из жизни. Турист прошёл \(\frac{2}{5}\) маршрута длиной \(15\) км. Сколько километров ему осталось?
\Needspace{3\baselineskip}\item[\textbf{3.16.}\starmark] Вычисли: \((\frac{3}{4} - \frac{5}{6}) : (\frac{7}{9} - \frac{2}{3}) + \frac{1}{2}\).
\Needspace{3\baselineskip}\item[\textbf{3.17.}\starmark] В первый день прочитали \(\frac{1}{3}\) книги, во второй — \(\frac{2}{5}\) оставшихся страниц. После этого осталось \(72\) страницы. Сколько страниц в книге?
\Needspace{3\baselineskip}\item[\textbf{3.18.}\starmark] Бак был полон. Израсходовали \(\frac{1}{4}\) воды, затем \(\frac{1}{3}\) оставшейся. После долива \(15\) л бак снова стал полным. Какова вместимость бака?
\end{enumerate}
\clearpage
\section{Десятичные дроби и проценты}
\subsection*{Краткая справка}
\begin{definition}\term{Десятичная дробь} — запись числа с разрядами десятых, сотых, тысячных и далее после запятой.\end{definition}
\begin{definition}\term{Процент} — одна сотая часть: \(1 \% = \frac{1}{100} = 0,01\).\end{definition}
\begin{definition}\term{Периодическая десятичная дробь} — бесконечная десятичная дробь с повторяющимся блоком цифр начиная с некоторого разряда.\end{definition}
\subsubsection*{Правила вычислений}
При сложении и вычитании десятичных дробей записывай одинаковые разряды друг под другом.\par
При умножении сначала умножь без запятых, затем отдели справа столько цифр, сколько их было после запятых в обоих множителях вместе.\par
При делении перенеси запятую в делимом и делителе на одинаковое число разрядов, чтобы делитель стал целым.\par
\subsubsection*{Формулы для процентов}
\(p \%\) от \(A\): \(A \cdot \frac{p}{100}\). Если \(p \%\) числа равны \(B\) и \(p \ne  0\), то число равно \(\frac{100 B}{p}\).\par
Рост на \(p \%\): \(A (1 + \frac{p}{100})\). Снижение на \(p \%\): \(A (1 - \frac{p}{100})\).\par
Процентное изменение: \(\frac{\text{новое значение} - \text{старое значение}}{\text{старое значение}} \cdot 100 \%\), если старое значение не равно \(0\). В обычных задачах о цене исходная цена положительна.\par
Каждый следующий процент изменения относится к уже изменённой величине.\par
\subsubsection*{Округление и периодические дроби}
Если следующая за оставляемым разрядом цифра \(0 - 4\), разряд не меняют; если \(5 - 9\), увеличивают на \(1\), при необходимости перенося единицу в старший разряд.\par
Любая периодическая десятичная дробь представляет рациональное число.\par
{\small\textit{Пример: \(\frac{1}{3} = 0{,}(3)\). Цена \(100\) рублей после двух скидок по \(10 \%\): \(100 \cdot 0,9 \cdot 0,9 = 81\) рубль.}}\par
\subsubsection*{Связь со следующими темами}
Десятичные приближения понадобятся при численном вычислении значений функций. Точное равенство и округлённый результат следует различать.
\Needspace{6\baselineskip}\subsection*{Задачи}
\begin{enumerate}[label=\thesection.\arabic*.,leftmargin=2.8em,itemsep=12pt]
\Needspace{3\baselineskip}\item Вычисли \(3,75 + 0,086\).
\Needspace{3\baselineskip}\item Вычисли \(8 - 2,47\).
\Needspace{3\baselineskip}\item Вычисли \(0,24 \cdot  0,5\).
\Needspace{3\baselineskip}\item Вычисли \(4,2 : 0,07\).
\Needspace{3\baselineskip}\item Запиши \(\frac{7}{8}\) десятичной дробью.
\Needspace{3\baselineskip}\item Запиши \(0,125\) несократимой обыкновенной дробью.
\Needspace{3\baselineskip}\item Запиши \(\frac{2}{3}\) периодической десятичной дробью.
\Needspace{3\baselineskip}\item Округли \(3,14159\) до сотых и до тысячных.
\Needspace{3\baselineskip}\item Найди \(15 \%\) от \(240\).
\Needspace{3\baselineskip}\item \(30 \%\) числа равны \(54\). Найди число.
\Needspace{3\baselineskip}\item Цена выросла с \(800\) до \(920\) рублей. На сколько процентов?
\Needspace{3\baselineskip}\item Цену \(1000\) рублей повысили на \(20 \%\), затем снизили на \(20 \%\). Найди итоговую цену и изменение в процентах от исходной.
\Needspace{3\baselineskip}\item Из жизни. Купили \(1,8\) кг яблок по \(125\) рублей за килограмм. Сколько заплатили?
\Needspace{3\baselineskip}\item Из жизни. Куртка стоит \(4800\) рублей. Скидка составляет \(15 \%\). Какова цена со скидкой?
\Needspace{3\baselineskip}\item Из жизни. Велосипедист проехал \(7,86\) км утром и \(5,47\) км вечером. Найди общий путь и округли его до десятых километра.
\Needspace{3\baselineskip}\item[\textbf{4.16.}\starmark] После скидки \(20 \%\) и следующей скидки \(10 \%\) товар стоит \(2160\) рублей. Найди исходную цену.
\Needspace{3\baselineskip}\item[\textbf{4.17.}\starmark] Цена выросла на \(25 \%\). На сколько процентов нужно снизить новую цену, чтобы получить исходную?
\Needspace{3\baselineskip}\item[\textbf{4.18.}\starmark] Смешали \(200\) г раствора с массовой долей соли \(10 \%\) и \(300\) г раствора с долей соли \(20 \%\). Найди массовую долю соли в смеси.
\end{enumerate}
\clearpage
\section{Отношения и пропорции}
\subsection*{Краткая справка}
\begin{definition}\term{Отношение} \(a\) к \(b\) — частное \(\frac{a}{b}\) при \(b \ne  0\); обозначается \(a : b\).\end{definition}
\begin{definition}\term{Пропорция} — равенство двух отношений: \(\frac{a}{b} = \frac{c}{d}\), где \(b,\allowbreak d \ne  0\).\end{definition}
\begin{definition}\term{Прямая пропорциональность} — зависимость \(y = kx\) с постоянным ненулевым коэффициентом \(k\).\end{definition}
\begin{definition}\term{Обратная пропорциональность} — зависимость \(y = \frac{k}{x}\) при \(x \ne  0\) с постоянным ненулевым \(k\).\end{definition}
\begin{definition}\term{Масштаб} \(1 : N\) — отношение длины на карте к длине на местности в одинаковых единицах.\end{definition}
\subsubsection*{Основное свойство пропорции}
При \(b,\allowbreak d \ne  0\) равенство \(\frac{a}{b} = \frac{c}{d}\) равносильно \(a \cdot  d = b \cdot  c\).\par
\subsubsection*{Правила применения}
При прямой пропорциональности увеличение \(x\) в несколько раз увеличивает \(y\) во столько же раз.\par
При обратной пропорциональности положительных величин увеличение \(x\) в несколько раз уменьшает \(y\) во столько же раз.\par
Деление \(S\) в отношении \(m : n\) при \(m,\allowbreak n > 0\): части равны \(S \cdot  \frac{m}{m + n}\) и \(S \cdot  \frac{n}{m + n}\).\par
{\small\textit{Пример: \(60\) разделить в отношении \(2 : 3\) — это \(24\) и \(36\), поскольку одна часть равна \(\frac{60}{5} = 12\).}}\par
\subsubsection*{Связь со следующими темами}
Прямая и обратная пропорциональность — первые примеры функций; их графики и свойства изучаются в главе 11.
\Needspace{6\baselineskip}\subsection*{Задачи}
\begin{enumerate}[label=\thesection.\arabic*.,leftmargin=2.8em,itemsep=12pt]
\Needspace{3\baselineskip}\item Сократи отношение \(18 : 24\).
\Needspace{3\baselineskip}\item Реши пропорцию \(\frac{x}{15} = \frac{4}{5}\).
\Needspace{3\baselineskip}\item Реши пропорцию \(\frac{7}{x} = \frac{21}{18}\).
\Needspace{3\baselineskip}\item Раздели \(64\) в отношении \(3 : 5\).
\Needspace{3\baselineskip}\item Для \(4\) порций нужно \(300\) г крупы. Сколько нужно для \(10\) порций?
\Needspace{3\baselineskip}\item \(5\) одинаковых тетрадей стоят \(180\) рублей. Сколько стоят \(8\)?
\Needspace{3\baselineskip}\item На карте масштаба \(1 : 100 000\) расстояние \(3,5\) см. Найди расстояние на местности в километрах.
\Needspace{3\baselineskip}\item \(6\) рабочих выполняют работу за \(10\) дней. За сколько дней справятся \(15\) рабочих при той же производительности?
\Needspace{3\baselineskip}\item Поездка при скорости \(60\) км/ч занимает \(3\) часа. Сколько займёт при \(90\) км/ч на том же пути?
\Needspace{3\baselineskip}\item Смешивают сок и воду в отношении \(2 : 3\). Сколько каждого компонента нужно для \(1,5\) л смеси?
\Needspace{3\baselineskip}\item При \(y = 3 x\) найди \(y\) для \(x = 4\). Как изменится \(y\) при удвоении \(x\)?
\Needspace{3\baselineskip}\item При \(y = \frac{24}{x}\) найди \(y\) для \(x = 6\). Как изменится \(y\) при удвоении положительного \(x\)?
\Needspace{3\baselineskip}\item Из жизни. Для \(6\) порций супа нужно \(450\) г картофеля. Сколько нужно для \(10\) таких же порций?
\Needspace{3\baselineskip}\item Из жизни. \(4\) одинаковых насоса заполняют бассейн за \(9\) часов. За сколько часов его заполнят \(6\) таких насосов, работающих одновременно?
\Needspace{3\baselineskip}\item Из жизни. Для напитка сироп смешивают с водой в отношении \(1 : 7\). Сколько сиропа и воды нужно для \(2\) л напитка?
\Needspace{3\baselineskip}\item[\textbf{5.16.}\starmark] Рабочие одинаковой производительности: \(6\) человек выполняют работу за \(8\) дней по \(5\) часов. За сколько дней справятся \(4\) человека по \(6\) часов?
\Needspace{3\baselineskip}\item[\textbf{5.17.}\starmark] В смеси отношение сиропа к воде \(2 : 5\). После добавления \(3\) л воды стало \(1 : 4\). Сколько литров смеси было сначала?
\Needspace{3\baselineskip}\item[\textbf{5.18.}\starmark] Автомобиль ехал одинаковое время со скоростями \(60\) и \(90\) км/ч. Найди среднюю скорость. А если одинаковыми были расстояния?
\end{enumerate}
\clearpage
\section{Степени с целым показателем}
\subsection*{Краткая справка}
\begin{definition}\term{Степень} \(a^{n}\) при натуральном \(n\) — произведение \(n\) одинаковых множителей \(a\).\end{definition}
\begin{definition}\term{Основание степени} \(a^{n}\) — число \(a\), которое возводят в степень.\end{definition}
\begin{definition}\term{Показатель степени} \(a^{n}\) — число \(n\); при натуральном \(n\) оно указывает количество одинаковых множителей \(a\).\end{definition}
\begin{definition}\term{Нулевая степень} ненулевого числа определяется равенством \(a^{0} = 1\).\end{definition}
\begin{definition}\term{Отрицательная целая степень} определяется равенством \(a^{- n} = \frac{1}{a^{n}}\) при \(a \ne  0\) и натуральном \(n\).\end{definition}
\begin{definition}\term{Стандартный вид положительного числа} — \(c \cdot  10^{n}\), где \(1 \le  c < 10\) и \(n\) — целое число.\end{definition}
\subsubsection*{Свойства степеней}
При умножении степеней одного основания показатели складывают: \(a^{m} \cdot  a^{n} = a^{m + n}\).\par
При делении показатели вычитают: \(\frac{a^{m}}{a^{n}} = a^{m - n},\allowbreak a \ne  0\).\par
При возведении степени в степень показатели перемножают: \((a^{m})^{n} = a^{mn}\).\par
Степень произведения: \((ab)^{n} = a^{n} b^{n}\). Степень дроби: \((\frac{a}{b})^{n} = \frac{a^{n}}{b^{n}},\allowbreak b \ne  0\).\par
Эти правила для целых показателей применяют только там, где все выражения определены; отрицательные степени требуют ненулевых оснований.\par
{\small\textit{Пример: \((- 2)^{4} = 16\), но \(- 2^{4} = - 16\). Число \(0,0032\) в стандартном виде: \(3,2 \cdot  10^{- 3}\).}}\par
\subsubsection*{Связь со следующими темами}
Буквенные правила степеней позволяют преобразовывать функции до вычисления пределов и производных.
\Needspace{6\baselineskip}\subsection*{Задачи}
\begin{enumerate}[label=\thesection.\arabic*.,leftmargin=2.8em,itemsep=12pt]
\Needspace{3\baselineskip}\item Вычисли \(2^{5}\).
\Needspace{3\baselineskip}\item Вычисли \((- 3)^{4}\) и \(- 3^{4}\).
\Needspace{3\baselineskip}\item Вычисли \(7^{0} + 2^{- 3}\).
\Needspace{3\baselineskip}\item Представь одной степенью: \(5^{3} \cdot  5^{4}\).
\Needspace{3\baselineskip}\item Представь одной степенью: \(2^{9} : 2^{4}\).
\Needspace{3\baselineskip}\item Представь одной степенью: \((3^{2})^{4}\).
\Needspace{3\baselineskip}\item Вычисли \((\frac{2}{3})^{- 2}\).
\Needspace{3\baselineskip}\item Вычисли \(10^{- 3} \cdot  10^{5}\).
\Needspace{3\baselineskip}\item Запиши \(0,00056\) в стандартном виде.
\Needspace{3\baselineskip}\item Запиши \(7,2 \cdot  10^{5}\) обычным числом.
\Needspace{3\baselineskip}\item Упрости \(a^{5} \cdot  a^{- 2}\) при \(a \ne  0\).
\Needspace{3\baselineskip}\item Вычисли \((- 2)^{3} \cdot  2^{- 2} + 3^{0}\).
\Needspace{3\baselineskip}\item Из жизни. В модели опыта каждая бактерия делится на две раз в час. Вначале была одна бактерия. Сколько их станет через \(8\) часов?
\Needspace{3\baselineskip}\item Из жизни. Квадратный участок имеет сторону \(10^{3}\) м. Запиши его площадь степенью десяти и обычным числом в м\(^{2}\).
\Needspace{3\baselineskip}\item Из жизни. В одном файле \(4 \cdot  10^{6}\) байт. Сколько байт в \(3 \cdot  10^{2}\) таких файлах?
\Needspace{3\baselineskip}\item[\textbf{6.16.}\starmark] Вычисли: \(\frac{2^{- 3} \cdot  4^{2}}{8^{- 1}}\).
\Needspace{3\baselineskip}\item[\textbf{6.17.}\starmark] Вычисли: \(\frac{3^{5} - 3^{4}}{3^{3} + 3^{2}}\).
\Needspace{3\baselineskip}\item[\textbf{6.18.}\starmark] В модели файл каждый день увеличивается втрое. Через \(4\) дня его размер \(162\) МБ. Каков был размер сначала и каким станет через \(6\) дней от начала?
\end{enumerate}
\clearpage
\section{Корни и действительные числа}
\subsection*{Краткая справка}
\begin{definition}\term{Арифметический квадратный корень} \(\sqrt{a}\) — неотрицательное число, квадрат которого равен \(a\). В действительных числах определён при \(a \ge  0\). При \(a>0\) уравнение \(z^2=a\) имеет два корня: \(z=\sqrt a\) и \(z=-\sqrt a\).\end{definition}
\begin{definition}\term{Кубический корень} \(\sqrt[3]{a}\) — число, куб которого равен \(a\); существует для любого действительного \(a\).\end{definition}
\begin{definition}\term{Рациональное число} — число вида \(\frac{p}{q}\), где \(p,\allowbreak q\) — целые и \(q \ne  0\).\end{definition}
\begin{definition}\term{Иррациональное число} — действительное число, которое нельзя представить в виде \(\frac{p}{q}\) с целыми \(p,q\) и \(q\ne0\).\end{definition}
\subsubsection*{Свойства корней}
При \(a,\allowbreak b \ge 0\): \(\sqrt{ab} = \sqrt{a} \cdot \sqrt{b}\).\par
При \(a \ge 0,\allowbreak b > 0\): \(\sqrt{\frac{a}{b}} = \frac{\sqrt{a}}{\sqrt{b}}\).\par
Для любого действительного \(x\): \(\sqrt{x^{2}} = | x |\).\par
Для \(a,\allowbreak b \ge  0\) из \(a < b\) следует \(\sqrt{a} < \sqrt{b}\).\par
{\small\textit{Примеры: \(\sqrt{72} = \sqrt{36 \cdot  2} = 6 \sqrt{2};\allowbreak \sqrt{(- 5)^{2}} = 5;\allowbreak \sqrt[3]{- 8} = - 2\).}}\par
\subsubsection*{Теорема Пифагора}
В прямоугольном треугольнике \(c^{2} = a^{2} + b^{2}\), где \(a,\allowbreak b\) — катеты, \(c\) — гипотенуза.\par
{\small\textit{Пример: при катетах \(3\) м и \(4\) м гипотенуза равна \(\sqrt{9 + 16} = 5\) м.}}\par
\subsubsection*{Связь со следующими темами}
В математическом анализе используются действительные числа: рациональные и иррациональные. Корень и его допустимые значения нужно различать.
\Needspace{6\baselineskip}\subsection*{Задачи}
\begin{enumerate}[label=\thesection.\arabic*.,leftmargin=2.8em,itemsep=12pt]
\Needspace{3\baselineskip}\item Вычисли \(\sqrt{81}\) и \(\sqrt[3]{- 27}\).
\Needspace{3\baselineskip}\item Вычисли \(\sqrt{\frac{49}{100}}\).
\Needspace{3\baselineskip}\item Упрости \(\sqrt{72}\).
\Needspace{3\baselineskip}\item Упрости \(\sqrt{200} - \sqrt{50}\).
\Needspace{3\baselineskip}\item Внеси множитель под корень: \(3 \sqrt{5}\).
\Needspace{3\baselineskip}\item Вычисли \(\sqrt{12} \cdot  \sqrt{3}\).
\Needspace{3\baselineskip}\item Вычисли \(\frac{\sqrt{75}}{\sqrt{3}}\).
\Needspace{3\baselineskip}\item Вычисли \(\sqrt{(- 7)^{2}}\).
\Needspace{3\baselineskip}\item Между какими соседними целыми числами находится \(\sqrt{30}\)?
\Needspace{3\baselineskip}\item Сравни \(\sqrt{7}\) и \(2,7\).
\Needspace{3\baselineskip}\item Какие числа рациональные: \(\sqrt{16},\allowbreak \sqrt{2},\allowbreak - \frac{3}{7},\allowbreak 0{,}(3),\allowbreak \pi \)?
\Needspace{3\baselineskip}\item Существуют ли \(\sqrt{- 9}\) и \(\sqrt[3]{- 8}\) в действительных числах?
\Needspace{3\baselineskip}\item Из жизни. Квадратная комната имеет площадь \(49\) м\(^{2}\). Найди длину её стороны.
\Needspace{3\baselineskip}\item Из жизни. Лестница стоит у стены: её основание в \(3\) м от стены, верх на высоте \(4\) м. Найди длину лестницы. Считай стену и землю перпендикулярными.
\Needspace{3\baselineskip}\item Из жизни. У кубического контейнера объём \(0,125\) м\(^{3}\). Найди длину ребра в метрах и сантиметрах.
\Needspace{3\baselineskip}\item[\textbf{7.16.}\starmark] Упрости: \(\sqrt{48} + \sqrt{75} - \sqrt{147}\).
\Needspace{3\baselineskip}\item[\textbf{7.17.}\starmark] Проверь без десятичных приближений утверждение \(\sqrt{2} + \sqrt{3} > \sqrt{10}\). Докажи его или опровергни.
\Needspace{3\baselineskip}\item[\textbf{7.18.}\starmark] У прямоугольника площадь \(48\) м\(^{2}\), одна сторона втрое длиннее другой. Найди обе стороны и диагональ.
\end{enumerate}
\clearpage
\section{Модуль и числовые промежутки}
\subsection*{Краткая справка}
\begin{definition}\term{Модуль} \(| x |\) — расстояние от числа \(x\) до нуля на числовой прямой: \(| x | = x\) при \(x \ge  0\) и \(| x | = - x\) при \(x < 0\).\end{definition}
\begin{definition}\term{Числовой промежуток} — множество всех чисел между заданными границами с указанным включением или исключением концов.\end{definition}
\begin{definition}\term{Пересечение} множеств \(A\) и \(B\), обозначаемое \(A\cap B\), — множество элементов, принадлежащих одновременно обоим множествам.\end{definition}
\begin{definition}\term{Объединение} множеств \(A\) и \(B\), обозначаемое \(A\cup B\), — множество элементов, принадлежащих хотя бы одному из этих множеств.\end{definition}
\subsubsection*{Свойства модуля}
\(| x | \ge  0;\allowbreak | x | = 0\) только при \(x = 0\). Также \(| - x | = | x |\) и \(| ab | = | a | | b |\).\par
Расстояние между числами \(a\) и \(b\) равно \(| a - b |\).\par
Неравенство треугольника: \(| a + b | \le  | a | + | b |\).\par
При \(r > 0\) условие \(| x - a | < r\) означает \(a - r < x < a + r\). При \(r>0\) условие \(|x-a|\le r\) равносильно \(a-r\le x\le a+r\).\par
\subsubsection*{Запись промежутков}
\((a;\allowbreak b)\): \(a < x < b\). \([a;\allowbreak b]\): \(a \le x \le b\). \([a;\allowbreak b)\) и \((a;\allowbreak b]\) включают только один из концов.\par
Бесконечные границы всегда записывают с круглой скобкой: например, \((- \infty ;\allowbreak 3]\).\par
{\small\textit{Пример: \(| x - 2 | \le  3\) равносильно \(- 1 \le  x \le  5\), то есть \(x \in  [- 1;\allowbreak 5]\).}}\par
\subsubsection*{Связь со следующими темами}
Запись через модуль описывает расстояние и близость чисел; на ней основаны определения предела и непрерывности.
\Needspace{6\baselineskip}\subsection*{Задачи}
\begin{enumerate}[label=\thesection.\arabic*.,leftmargin=2.8em,itemsep=12pt]
\Needspace{3\baselineskip}\item Вычисли \(| - 8 | + | 3 |\).
\Needspace{3\baselineskip}\item Вычисли \(| - 5 + 2 |\).
\Needspace{3\baselineskip}\item Найди расстояние между \(- 7\) и \(4\).
\Needspace{3\baselineskip}\item Сравни \(| - 6 |\) и \(| 5 |\).
\Needspace{3\baselineskip}\item Запиши \(- 2 < x \le  3\) промежутком.
\Needspace{3\baselineskip}\item Запиши \([- 4;\allowbreak 1)\) двойным неравенством.
\Needspace{3\baselineskip}\item Выпиши целые числа из \((- 3;\allowbreak 2]\).
\Needspace{3\baselineskip}\item Найди пересечение \([- 2;\allowbreak 4]\) и \((1;\allowbreak 7)\).
\Needspace{3\baselineskip}\item Найди объединение \((- \infty ;\allowbreak 2]\) и \((1;\allowbreak 5)\).
\Needspace{3\baselineskip}\item Реши \(| x | = 5\).
\Needspace{3\baselineskip}\item Реши \(| x - 2 | < 3\).
\Needspace{3\baselineskip}\item Упрости \(\sqrt{x^{2}}\) и вычисли при \(x = - 4\).
\Needspace{3\baselineskip}\item Из жизни. Температура в морозилке \(- 18 {}^{\circ} C\), а в комнате \(+ 22 {}^{\circ} C\). На сколько градусов отличаются температуры?
\Needspace{3\baselineskip}\item Из жизни. Длина детали должна отличаться от \(50\) мм не более чем на \(0,2\) мм. Запиши допустимые длины \(x\) через модуль и промежутком.
\Needspace{3\baselineskip}\item Из жизни. Аттракцион доступен при росте от \(120\) см включительно до \(190\) см не включительно. Запиши условие для роста \(h\) неравенством и промежутком. Можно ли войти при росте \(190\) см?
\Needspace{3\baselineskip}\item[\textbf{8.16.}\starmark] Реши \(| x - 3 | + | x + 1 | = 6\).
\Needspace{3\baselineskip}\item[\textbf{8.17.}\starmark] Найди все целые \(x\), для которых \(| 2 x - 1 | \le  5\) и \(x \notin  [0;\allowbreak 2)\).
\Needspace{3\baselineskip}\item[\textbf{8.18.}\starmark] Два датчика допускают температуру: первый \(| t - 20 | \le  3\), второй \(| t - 24 | < 4\). Найди общий допустимый диапазон.
\end{enumerate}
\clearpage
\section{Преобразование алгебраических выражений}
\subsection*{Краткая справка}
\begin{definition}\term{Подобные слагаемые} — слагаемые с одинаковой буквенной частью, включая показатели степеней.\end{definition}
\begin{definition}\term{Разложение на множители} — представление выражения в виде произведения.\end{definition}
\begin{definition}\term{Тождество} — равенство, верное при всех допустимых значениях переменных.\end{definition}
\begin{definition}\term{Область допустимых значений (ОДЗ)} — значения переменных, при которых исходное выражение имеет смысл.\end{definition}
\subsubsection*{Формулы сокращённого умножения}
\((a + b)^{2} = a^{2} + 2 ab + b^{2}\).\par
\((a - b)^{2} = a^{2} - 2 ab + b^{2}\).\par
\(a^{2} - b^{2} = (a - b) (a + b)\).\par
\subsubsection*{Правила преобразований}
Общий множитель: \(ab + ac = a (b + c)\). Подобные слагаемые складывают по коэффициентам: \(2 x + 3 x = 5 x\).\par
В дроби сокращают общие множители числителя и знаменателя, а не отдельные слагаемые.\par
После сокращения сохраняют запреты исходного выражения. Знаменатель не равен \(0\); под квадратным корнем стоит неотрицательное число.\par
Чтобы убрать корень из знаменателя вида \(\sqrt{a} - b\), можно умножить числитель и знаменатель на \(\sqrt{a} + b\), если преобразование допустимо.\par
{\small\textit{Пример: \(\frac{x^{2} - 1}{x - 1} = x + 1\), но только при \(x \ne  1\).}}\par
\subsubsection*{Связь со следующими темами}
Перед вычислением предела часто требуется преобразовать выражение. Равенство формул на общей ОДЗ не отменяет запретов исходной записи.
\Needspace{6\baselineskip}\subsection*{Задачи}
\begin{enumerate}[label=\thesection.\arabic*.,leftmargin=2.8em,itemsep=12pt]
\Needspace{3\baselineskip}\item Раскрой скобки: \(3 (2 x - 5) - 2 (x + 1)\).
\Needspace{3\baselineskip}\item Приведи подобные: \(4 a - 3 b + 2 a + 5 b\).
\Needspace{3\baselineskip}\item Раскрой скобки: \((x + 3)^{2}\).
\Needspace{3\baselineskip}\item Раскрой скобки: \((2 x - 1) (2 x + 1)\).
\Needspace{3\baselineskip}\item Разложи на множители: \(6 x^{2} - 9 x\).
\Needspace{3\baselineskip}\item Разложи на множители: \(x^{2} - 16\).
\Needspace{3\baselineskip}\item Разложи на множители: \(x^{2} + 5 x + 6\).
\Needspace{3\baselineskip}\item Сократи \(\frac{x^{2} - 1}{x - 1}\). Укажи ограничения исходной дроби.
\Needspace{3\baselineskip}\item Упрости \(\frac{2}{x} + \frac{3}{2 x}\). Укажи ограничения.
\Needspace{3\baselineskip}\item Упрости \(\frac{a^{2} b}{ab^{2}}\). Укажи ограничения.
\Needspace{3\baselineskip}\item Упрости \(\frac{x^{2} - 4}{x^{2} + 4 x + 4}\). Укажи ограничения.
\Needspace{3\baselineskip}\item Упрости \(\frac{1}{\sqrt{5} - 2}\), избавившись от корня в знаменателе.
\Needspace{3\baselineskip}\item Из жизни. Билет стоит \(x\) рублей. Заказали \(3\) взрослых билета по полной цене и \(2\) детских по цене \(x - 100\) рублей. Составь и упрости выражение общей стоимости, если \(x > 100\).
\Needspace{3\baselineskip}\item Из жизни. У прямоугольной клумбы стороны \(x + 2\) и \(x + 5\) метров, где \(x > 0\). Вырази площадь многочленом.
\Needspace{3\baselineskip}\item Из жизни. Квадратную площадку со стороной \(x>0\) метров расширили на \(3\) м в каждом измерении. На сколько выросла её площадь?
\Needspace{3\baselineskip}\item[\textbf{9.16.}\starmark] Упрости \(\frac{x^{2} - 9}{x^{2} - x - 6}\) и укажи ОДЗ исходного выражения.
\Needspace{3\baselineskip}\item[\textbf{9.17.}\starmark] Упрости: \((\frac{1}{x - 1} - \frac{1}{x + 1}) : (\frac{2}{x^{2} - 1})\).
\Needspace{3\baselineskip}\item[\textbf{9.18.}\starmark] Докажи тождество \((a + b)^{2} - (a - b)^{2} = 4 ab\) и проверь его при \(a = 103,\allowbreak b = 97\), не вычисляя квадраты больших чисел.
\end{enumerate}
\clearpage
\section{Уравнения, системы и неравенства}
\subsection*{Краткая справка}
\begin{definition}\term{Корень уравнения} — допустимое значение неизвестного, при котором равенство верно.\end{definition}
\begin{definition}\term{Решение неравенства} — значение неизвестного, при котором указанное сравнение верно.\end{definition}
\begin{definition}\term{Решение системы} — значения неизвестных, которые удовлетворяют всем её условиям одновременно.\end{definition}
\subsubsection*{Линейные и квадратные уравнения}
\(ax + b = 0\) при \(a \ne  0\) имеет корень \(x = - \frac{b}{a}\). При \(a = 0,\allowbreak b = 0\) подходят все \(x\); при \(a = 0,\allowbreak b \ne  0\) решений нет.\par
\begin{definition}\term{Дискриминант} квадратного уравнения \(ax^{2}+bx+c=0\) при \(a\ne0\) — число \(D=b^{2}-4ac\).\end{definition}
При \(D > 0\) два корня: \(x = \frac{- b \pm  \sqrt{D}}{2 a}\). При \(D = 0\) один корень: \(- \frac{b}{2 a}\). При \(D < 0\) действительных корней нет.\par
\subsubsection*{Правила и проверка}
К обеим частям равенства или неравенства можно прибавить одно и то же выражение при сохранении ОДЗ.\par
При умножении или делении неравенства на отрицательное число знак меняется; на положительное — сохраняется.\par
Метод интервалов: отметь нули и запрещённые точки, определи знак на каждом участке и отдельно проверь включение границ.\par
При возведении уравнения в квадрат могут появиться посторонние корни: проверь найденные значения в исходном условии.\par
{\small\textit{Пример: \(- 2 x > 6\) даёт \(x < - 3\). Уравнение \(\sqrt{x + 2} = x\) требует \(x \ge  0\) ещё до возведения в квадрат.}}\par
\subsubsection*{Связь со следующими темами}
Множество решений — это все подходящие значения, а не один найденный пример. Такой подход сохраняется и в университетских задачах.
\Needspace{6\baselineskip}\subsection*{Задачи}
\begin{enumerate}[label=\thesection.\arabic*.,leftmargin=2.8em,itemsep=12pt]
\Needspace{3\baselineskip}\item Реши \(3 x - 7 = 11\).
\Needspace{3\baselineskip}\item Реши \(2 (x - 3) = 5 x + 9\).
\Needspace{3\baselineskip}\item Реши \(x^{2} - 5 x + 6 = 0\).
\Needspace{3\baselineskip}\item Реши \(2 x^{2} - 8 = 0\).
\Needspace{3\baselineskip}\item Реши \(\frac{x + 1}{x - 2} = 2\).
\Needspace{3\baselineskip}\item Реши систему: \(x + y = 7;\allowbreak x - y = 1\).
\Needspace{3\baselineskip}\item Реши \(- 3 x > 12\).
\Needspace{3\baselineskip}\item Реши \(x^{2} - 9 \le  0\).
\Needspace{3\baselineskip}\item Реши \((x - 1) (x + 2) > 0\).
\Needspace{3\baselineskip}\item Реши \(\frac{x - 3}{x + 1} \le  0\).
\Needspace{3\baselineskip}\item Реши \(| 2 x - 1 | = 5\).
\Needspace{3\baselineskip}\item Реши \(\sqrt{x + 2} = x\).
\Needspace{3\baselineskip}\item Из жизни. Поездка на такси стоит \(150\) рублей за подачу и \(30\) рублей за километр. Счёт составил \(510\) рублей. Какое расстояние проехали?
\Needspace{3\baselineskip}\item Из жизни. За \(3\) одинаковых блокнота и ручку стоимостью \(80\) рублей заплатили \(440\) рублей. Сколько стоит блокнот?
\Needspace{3\baselineskip}\item Из жизни. Для покупки билетов по \(350\) рублей выделено \(2000\) рублей. Какое максимальное целое число билетов можно купить?
\Needspace{3\baselineskip}\item[\textbf{10.16.}\starmark] Реши \(\sqrt{2 x + 3} = x\).
\Needspace{3\baselineskip}\item[\textbf{10.17.}\starmark] Реши \(\frac{x - 2}{x + 3} \ge  1\).
\Needspace{3\baselineskip}\item[\textbf{10.18.}\starmark] У прямоугольного участка периметр \(34\) м, площадь \(60\) м\(^{2}\). Найди стороны.
\end{enumerate}
\clearpage
\section{Функции и графики}
\subsection*{Краткая справка}
\begin{definition}\term{Функция} — соответствие, которое каждому \(x\) из области определения ставит в соответствие ровно одно значение \(f (x)\).\end{definition}
\begin{definition}\term{Область определения} функции — множество всех допустимых значений аргумента \(x\).\end{definition}
\begin{definition}\term{Область значений} функции — множество всех значений \(f(x)\), которые она принимает.\end{definition}
\begin{definition}\term{График функции} \(f\) — множество точек \((x;f(x))\) при всех \(x\) из её области определения.\end{definition}
\begin{definition}\term{Нули функции} — значения \(x\) из области определения, для которых \(f(x)=0\).\end{definition}
\begin{definition}\term{Строгое возрастание} функции \(f\) на промежутке \(I\) означает: для любых \(x_1,x_2\in I\) из \(x_1<x_2\) следует \(f(x_1)<f(x_2)\).\end{definition}
\begin{definition}\term{Строгое убывание} функции \(f\) на промежутке \(I\) означает: для любых \(x_1,x_2\in I\) из \(x_1<x_2\) следует \(f(x_1)>f(x_2)\).\end{definition}
В этом сборнике слова «возрастает» и «убывает» означают строгую монотонность.
\begin{definition}\term{Неубывающая функция} на промежутке \(I\) — функция, для которой при любых \(x_1,x_2\in I\) из \(x_1<x_2\) следует \(f(x_1)\le f(x_2)\). Это более слабое условие, чем строгое возрастание.\end{definition}
\begin{definition}\term{Чётная функция} — функция с симметричной относительно нуля областью определения, для которой \(f(-x)=f(x)\) при всех допустимых \(x\). Симметричность области означает, что вместе с \(x\) она содержит \(-x\).\end{definition}
\begin{definition}\term{Нечётная функция} — функция с симметричной относительно нуля областью определения, для которой \(f(-x)=-f(x)\) при всех допустимых \(x\).\end{definition}
\subsubsection*{Графики и преобразования}
\(y = f (x - a)\) при \(a > 0\) сдвигает график вправо на \(a\). \(y = f (x) + b\) при \(b > 0\) сдвигает его вверх на \(b\).\par
\(y = - f (x)\) отражает график относительно горизонтальной оси. \(y = f (- x)\) отражает его относительно вертикальной оси.\par
Чтобы найти пересечения двух графиков, реши \(f (x) = g (x)\), затем вычисли \(y\) для каждого найденного \(x\).\par
\begin{definition}\term{Композиция} функций \(f\) и \(g\), обозначаемая \(f(g(x))\), — функция, для вычисления которой сначала находят \(g(x)\), затем подставляют результат в \(f\). Аргумент \(x\) должен принадлежать области определения \(g\), а значение \(g(x)\) — области определения \(f\).\end{definition}
{\small\textit{Пример: \(y = (x - 2)^{2} + 1\) имеет вершину \((2;\allowbreak 1)\) и область значений \([1;\allowbreak + \infty )\).}}\par
\subsubsection*{Связь со следующими темами}
В математическом анализе изучают изменение функций. Поэтому до производных важно понимать аргумент, область определения, график и монотонность.
\Needspace{6\baselineskip}\subsection*{Задачи}
\begin{enumerate}[label=\thesection.\arabic*.,leftmargin=2.8em,itemsep=12pt]
\Needspace{3\baselineskip}\item Для \(f (x) = 2 x - 3\) найди \(f (- 2),\allowbreak f (0),\allowbreak f (3)\).
\Needspace{3\baselineskip}\item Найди область определения \(f (x) = \frac{1}{x - 4}\).
\Needspace{3\baselineskip}\item Найди область определения \(f (x) = \sqrt{x + 3}\).
\Needspace{3\baselineskip}\item Найди нули \(f (x) = x^{2} - 4\).
\Needspace{3\baselineskip}\item Где \(f (x) = 3 x - 6\) положительна?
\Needspace{3\baselineskip}\item Построй \(y = 2 x + 1\) и найди пересечения с осями.
\Needspace{3\baselineskip}\item Построй \(y = (x - 2)^{2} - 1\). Укажи вершину.
\Needspace{3\baselineskip}\item Укажи промежутки убывания и возрастания \(y = x^{2}\).
\Needspace{3\baselineskip}\item Определи чётность \(f (x) = x^{3}\) и \(g (x) = x^{2} + 1\).
\Needspace{3\baselineskip}\item При \(f (x) = x^{2}\) и \(g (x) = x + 1\) найди \(f (g (x))\) и \(g (f (x))\).
\Needspace{3\baselineskip}\item Найди область значений \(y = (x - 1)^{2} + 2\).
\Needspace{3\baselineskip}\item Найди точки пересечения \(y = x^{2}\) и \(y = x + 2\).
\Needspace{3\baselineskip}\item Из жизни. Стоимость аренды велосипеда описана формулой \(C (t) = 200 + 150 t\) рублей, где \(t \ge  0\) — время в часах. Найди стоимость для \(0,\allowbreak 2\) и \(4\) часов и построй график.
\Needspace{3\baselineskip}\item Из жизни. В баке \(60\) л воды. Её расходуют со скоростью \(5\) л в минуту, не доливая. Запиши объём \(V (t)\), укажи допустимое время и момент опустошения.
\Needspace{3\baselineskip}\item Из жизни. Два тарифа доставки: \(A (x) = 100 + 20 x\) и \(B (x) = 40 x\) рублей, где \(x \ge  0\) — расстояние в километрах. При каком расстоянии цены равны? Какой тариф дешевле при \(2\) км и при \(8\) км?
\Needspace{3\baselineskip}\item[\textbf{11.16.}\starmark] Найди область определения \(f (x) = \frac{\sqrt{x - 1}}{x - 3}\).
\Needspace{3\baselineskip}\item[\textbf{11.17.}\starmark] Построй \(y = | x - 2 | + 1\) и найди все \(x\), для которых \(y \le  4\).
\Needspace{3\baselineskip}\item[\textbf{11.18.}\starmark] Из куска проволоки длиной \(20\) м делают прямоугольник. Вырази площадь через сторону \(x\) и найди наибольшую площадь.
\end{enumerate}
\clearpage
\section{Степенные, показательные и логарифмические функции}
\subsection*{Краткая справка}
\begin{definition}\term{Степенная функция} — функция \(y = x^{r}\) с фиксированным показателем \(r\); допустимые \(x\) зависят от показателя.\end{definition}
\begin{definition}\term{Показательная функция} — \(y = a^{x}\) при \(a > 0\) и \(a \ne  1;\allowbreak x\) может быть любым действительным числом.\end{definition}
\begin{definition}\term{Логарифм} \(\log _{a} b\) — показатель степени, в которую нужно возвести \(a\), чтобы получить \(b\): \(a^{c} = b\) означает \(\log _{a} b = c\).\end{definition}
Условия логарифма: \(a > 0,\allowbreak a \ne  1,\allowbreak b > 0\).\par
\subsubsection*{Свойства показательной и логарифмической функций}
\(a^{x}\) всегда положительно. При \(a > 1\) функция возрастает, при \(0 < a < 1\) убывает. Логарифм с тем же основанием меняется в том же направлении.\par
\(\log _{a} 1 = 0;\allowbreak \log _{a} a = 1\). При допустимых значениях \(\log _{a} (a^{x}) = x\) и \(a^{\log _{a} b} = b\).\par
При \(u,\allowbreak v > 0\): \(\log _{a} (uv) = \log _{a} u + \log _{a} v;\allowbreak \log _{a} (\frac{u}{v}) = \log _{a} u - \log _{a} v\).\par
При \(u > 0\): \(\log _{a} (u^{r}) = r \log _{a} u\). Формула смены основания: \(\log _{a} u = \frac{\log _{c} u}{\log _{c} a}\) при \(c > 0,\allowbreak c \ne  1\).\par
Графики \(y = a^{x}\) и \(y = \log _{a} x\) симметричны относительно прямой \(y = x\).\par
\subsubsection*{Две важные степенные функции}
\(y = \frac{1}{x}\): \(x \ne 0\), значения \(y \ne  0\). \(y = \sqrt{x}\): \(x \ge 0\), значения \(y \ge  0\).\par
{\small\textit{Пример: \(\log _{2} 8 = 3\), потому что \(2^{3} = 8\). Неравенство \((\frac{1}{2})^{x} > (\frac{1}{2})^{3}\) даёт \(x < 3\).}}\par
\subsubsection*{Связь со следующими темами}
Степенные, показательные и логарифмические функции образуют базовый набор функций для изучения пределов, производных и интегралов.
\Needspace{6\baselineskip}\subsection*{Задачи}
\begin{enumerate}[label=\thesection.\arabic*.,leftmargin=2.8em,itemsep=12pt]
\Needspace{3\baselineskip}\item Построй \(y = \frac{1}{x}\). Укажи область определения и область значений.
\Needspace{3\baselineskip}\item Построй \(y = \sqrt{x}\). Укажи область определения и область значений.
\Needspace{3\baselineskip}\item Реши \(x^{3} = - 27\).
\Needspace{3\baselineskip}\item Реши \(2^{x} = 32\).
\Needspace{3\baselineskip}\item Реши \(3^{x - 1} = 9\).
\Needspace{3\baselineskip}\item Реши \((\frac{1}{2})^{x} > \frac{1}{8}\).
\Needspace{3\baselineskip}\item Вычисли \(\log _{2} 8 + \log _{3} 1\).
\Needspace{3\baselineskip}\item Вычисли \(\log _{5} (\frac{1}{25})\).
\Needspace{3\baselineskip}\item Реши \(\log _{2} (x - 1) = 3\).
\Needspace{3\baselineskip}\item Найди область определения \(y = \log _{3} (5 - x)\).
\Needspace{3\baselineskip}\item При \(a > 0\) упрости \(\log _{2} (8 a) - \log _{2} a\).
\Needspace{3\baselineskip}\item Сравни \(2^{0,3}\) и \(2^{0,7}\); затем \((\frac{1}{3})^{0,3}\) и \((\frac{1}{3})^{0,7}\).
\Needspace{3\baselineskip}\item Из жизни. В учебной модели стоимость оборудования каждый год уменьшается вдвое: \(S (t) = 80000 \cdot  (\frac{1}{2})^{t}\) рублей. Найди стоимость через \(3\) года.
\Needspace{3\baselineskip}\item Из жизни. Модель роста культуры: \(N (t) = 500 \cdot  2^{t}\), где \(t\) — число часов. Через сколько часов станет \(8000\) клеток?
\Needspace{3\baselineskip}\item Из жизни. Для двоичного поиска в массиве из \(2^{k}\) элементов используют модель из \(k\) шагов. Сколько шагов даёт эта модель для \(1024\) элементов?
\Needspace{3\baselineskip}\item[\textbf{12.16.}\starmark] Реши \(4^{x} - 5 \cdot  2^{x} + 4 = 0\).
\Needspace{3\baselineskip}\item[\textbf{12.17.}\starmark] Реши \(\log _{2} (x - 1) + \log _{2} (x + 1) = 3\).
\Needspace{3\baselineskip}\item[\textbf{12.18.}\starmark] Вклад в учебной модели ежегодно увеличивается на \(10 \%\) без пополнений и снятий. Найди наименьшее целое число лет, через которое \(10000\) рублей станут не меньше \(13310\) рублей.
\end{enumerate}
\clearpage
\section{Тригонометрия}
\subsection*{Краткая справка}
\begin{definition}\term{Радиан} — единица измерения угла. Угол в \(1\) радиан — центральный угол, которому соответствует дуга длины, равной радиусу окружности. Радианная мера центрального угла равна \(\alpha=\frac{l}{r}\), где \(l\) — длина дуги, \(r>0\) — радиус. Полный оборот: \(2\pi\) радиан, то есть \(360^\circ\).\end{definition}
Положительные углы отсчитывают от точки \((1;0)\) единичной окружности против часовой стрелки, отрицательные — по часовой.
\begin{definition}\term{Косинус} \(\cos\alpha\) — горизонтальная координата точки единичной окружности, полученной поворотом точки \((1;0)\) на угол \(\alpha\).\end{definition}
\begin{definition}\term{Синус} \(\sin\alpha\) — вертикальная координата точки единичной окружности, полученной поворотом точки \((1;0)\) на угол \(\alpha\).\end{definition}
\begin{definition}\term{Тангенс}: \(\tan  \alpha  = \frac{\sin  \alpha }{\cos  \alpha }\) при \(\cos  \alpha  \ne  0\).\end{definition}
\begin{definition}\term{Период функции} — ненулевое число \(T\), при сдвиге аргумента на которое область определения сохраняется и \(f (x + T) = f (x)\).\end{definition}
\subsubsection*{Основные тождества и свойства}
\(\sin ^{2} \alpha  + \cos ^{2} \alpha  = 1\). Значения синуса и косинуса лежат в \([- 1;\allowbreak 1]\).\par
\(\sin  (- \alpha ) = - \sin  \alpha ;\allowbreak \cos  (- \alpha ) = \cos  \alpha \).\par
\(\sin  (\pi  - \alpha ) = \sin  \alpha ;\allowbreak \cos  (\pi  - \alpha ) = - \cos  \alpha \).\par
Наименьший положительный период синуса и косинуса равен \(2 \pi \). Для \(\sin  (kx)\) и \(\cos  (kx)\) при \(k \ne  0\) он равен \(\frac{2 \pi }{| k |}\).\par
\subsubsection*{Прямоугольный треугольник}
Для острого угла синус — \(\frac{\text{противолежащий катет}}{\text{гипотенуза}}\); косинус — \(\frac{\text{прилежащий катет}}{\text{гипотенуза}}\); тангенс — \(\frac{\text{противолежащий катет}}{\text{прилежащий катет}}\).\par
\subsubsection*{Значения, которые стоит знать}
\(\sin  0 = 0;\allowbreak \sin  (\frac{\pi }{6}) = \frac{1}{2};\allowbreak \sin  (\frac{\pi }{4}) = \frac{\sqrt{2}}{2};\allowbreak \sin  (\frac{\pi }{3}) = \frac{\sqrt{3}}{2};\allowbreak \sin  (\frac{\pi }{2}) = 1\).\par
\(\cos  0 = 1;\allowbreak \cos  (\frac{\pi }{6}) = \frac{\sqrt{3}}{2};\allowbreak \cos  (\frac{\pi }{4}) = \frac{\sqrt{2}}{2};\allowbreak \cos  (\frac{\pi }{3}) = \frac{1}{2};\allowbreak \cos  (\frac{\pi }{2}) = 0\).\par
{\small\textit{Пример: \(60 {}^{\circ} = \frac{\pi }{3}\) радиан, потому что \(60 \cdot  \frac{\pi }{180} = \frac{\pi }{3}\).}}\par
\subsubsection*{Связь со следующими темами}
Радианная мера необходима для стандартных формул пределов и производных тригонометрических функций.
\Needspace{6\baselineskip}\subsection*{Задачи}
\begin{enumerate}[label=\thesection.\arabic*.,leftmargin=2.8em,itemsep=12pt]
\Needspace{3\baselineskip}\item Переведи \(30 {}^{\circ},\allowbreak 90 {}^{\circ}\) и \(150 {}^{\circ}\) в радианы.
\Needspace{3\baselineskip}\item Переведи \(\frac{\pi }{4}\) и \(\frac{7 \pi }{6}\) в градусы.
\Needspace{3\baselineskip}\item Вычисли \(\sin  (\frac{\pi }{6}),\allowbreak \cos  (\frac{\pi }{3}),\allowbreak \tan  (\frac{\pi }{4})\).
\Needspace{3\baselineskip}\item Определи знаки \(\sin \) и \(\cos \) для угла \(\frac{5 \pi }{6}\).
\Needspace{3\baselineskip}\item Найди \(\sin  (- \frac{\pi }{6})\) и \(\cos  (- \frac{\pi }{3})\).
\Needspace{3\baselineskip}\item Упрости \(\sin ^{2} x + \cos ^{2} x + 2\).
\Needspace{3\baselineskip}\item Если \(\sin  x = \frac{3}{5}\) и \(0 < x < \frac{\pi }{2}\), найди \(\cos  x\) и \(\tan  x\).
\Needspace{3\baselineskip}\item Найди \(\sin  (\pi  - x)\) и \(\cos  (\pi  - x)\) через функции \(x\).
\Needspace{3\baselineskip}\item Реши \(\sin  x = 0\) при \(x \in  [0;\allowbreak 2 \pi ]\).
\Needspace{3\baselineskip}\item Реши \(\cos  x = \frac{1}{2}\) при \(x \in  [0;\allowbreak 2 \pi ]\).
\Needspace{3\baselineskip}\item Построй \(y = \sin  x\) на \([0;\allowbreak 2 \pi ]\). Укажи нули, максимум и минимум.
\Needspace{3\baselineskip}\item Найди наименьший положительный период \(y = \sin  (2 x)\).
\Needspace{3\baselineskip}\item Из жизни. Солнечный луч образует с землёй угол \(45 {}^{\circ}\). Тень вертикального столба равна \(6\) м. Найди высоту столба.
\Needspace{3\baselineskip}\item Из жизни. Лестница длиной \(10\) м образует с горизонтальной землёй угол \(30 {}^{\circ}\). На какой высоте находится её верх?
\Needspace{3\baselineskip}\item Из жизни. Колесо обозрения радиусом \(10\) м имеет центр на высоте \(12\) м. Высота кабинки \(h = 12 + 10 \sin  \alpha \). Найди высоту при \(\alpha  = 0,\allowbreak \frac{\pi }{2}\) и \(\frac{3 \pi }{2}\).
\Needspace{3\baselineskip}\item[\textbf{13.16.}\starmark] Реши \(2 \sin ^{2} x - 3 \sin  x + 1 = 0\) при \(x \in  [0;\allowbreak 2 \pi ]\).
\Needspace{3\baselineskip}\item[\textbf{13.17.}\starmark] Луч образует с землёй угол \(30 {}^{\circ}\). Длина тени вертикального дерева \(12\) м. Найди высоту дерева в точном виде.
\Needspace{3\baselineskip}\item[\textbf{13.18.}\starmark] Высота кабинки задана \(h (t) = 15 - 12 \cos  (\frac{\pi  t}{20})\) метров, где \(t\) — секунды от начала движения. Найди минимальную и максимальную высоты и время первого достижения максимума при \(t \ge  0\).
\end{enumerate}
\clearpage
\section{Множества, логика и доказательства}
\subsection*{Краткая справка}
\begin{definition}\term{Множество} — совокупность элементов. Запись \(x \in  A\) означает, что \(x\) принадлежит множеству \(A\).\end{definition}
\begin{definition}\term{Пересечение} \(A\cap B\) — множество элементов, принадлежащих одновременно \(A\) и \(B\).\end{definition}
\begin{definition}\term{Объединение} \(A\cup B\) — множество элементов, принадлежащих хотя бы одному из множеств \(A\) и \(B\).\end{definition}
\begin{definition}\term{Утверждение} — высказывание, для которого можно определить истинность или ложность в заданных условиях.\end{definition}
\begin{definition}\term{Импликация} \(P\Rightarrow Q\) — высказывание «если \(P\), то \(Q\)».\end{definition}
\begin{definition}\term{Достаточное условие} для \(Q\) — условие \(P\), из которого следует \(Q\), то есть импликация \(P\Rightarrow Q\) истинна.\end{definition}
\begin{definition}\term{Необходимое условие} для \(P\) — условие \(Q\), которое обязательно выполняется при выполнении \(P\), то есть импликация \(P\Rightarrow Q\) истинна.\end{definition}
\subsubsection*{Кванторы и отрицание}
\(\forall \) означает «для любого», \(\exists \) — «существует». \(\mathbb{Z}\) обозначает целые числа, \(\mathbb{R}\) — действительные.\par
Отрицание «для всех \(x\) верно \(P (x)\)» — «существует \(x\), для которого \(P (x)\) неверно».\par
Чтобы опровергнуть всеобщее утверждение, достаточно одного контрпримера. Чтобы доказать его, проверки нескольких примеров недостаточно.\par
\subsubsection*{Свойства и способы доказательства}
Для конечных множеств: \(| A \cup  B | = | A | + | B | - | A \cap  B |\). Здесь \(| A |\) означает число элементов множества, а не модуль числа.\par
\begin{definition}\term{Контрапозиция} импликации \(P\Rightarrow Q\) — импликация \(\neg Q\Rightarrow\neg P\), равносильная исходной:\[ (P\Rightarrow Q)\Longleftrightarrow(\neg Q\Rightarrow\neg P). \] Здесь \(\neg\) означает отрицание.\end{definition}
Обратное утверждение \(Q\Rightarrow P\) в общем случае из \(P\Rightarrow Q\) не следует.\par
Принцип математической индукции: проверь утверждение для \(n = 1\); затем для любого \(k \ge  1\) докажи переход от его истинности при \(n = k\) к истинности при \(n = k + 1\). Тогда оно верно для всех натуральных \(n \ge  1\).\par
{\small\textit{Пример прямого доказательства: сумма чётных чисел \(2 m\) и \(2 n\) равна \(2 (m + n)\), поэтому она чётна.}}\par
\subsubsection*{Связь со следующими темами}
В определениях предела важен порядок слов «для любого» и «существует». Различайте само утверждение, его отрицание и обратное утверждение.
\Needspace{6\baselineskip}\subsection*{Задачи}
\begin{enumerate}[label=\thesection.\arabic*.,leftmargin=2.8em,itemsep=12pt]
\Needspace{3\baselineskip}\item При \(A = \{1,\allowbreak 2,\allowbreak 3\}\) и \(B = \{3,\allowbreak 4\}\) найди \(A \cap  B\) и \(A \cup  B\).
\Needspace{3\baselineskip}\item Запиши символами: для любого действительного \(x\) число \(x^{2}\) неотрицательно.
\Needspace{3\baselineskip}\item Переведи: \(\exists  x \in  \mathbb{Z} : x^{2} = 9\). Верно ли это?
\Needspace{3\baselineskip}\item Верно ли: для любого действительного \(x\) выполнено \(x^{2} > 0\)?
\Needspace{3\baselineskip}\item Запиши отрицание фразы «все числа в наборе положительны».
\Needspace{3\baselineskip}\item Верна ли импликация: если число делится на \(4\), то оно чётное? Верно ли обратное?
\Needspace{3\baselineskip}\item Для делимости на \(6\) чётность — необходимое или достаточное условие?
\Needspace{3\baselineskip}\item Докажи, что сумма двух чётных целых чисел чётна.
\Needspace{3\baselineskip}\item Докажи, что квадрат нечётного целого числа нечётный.
\Needspace{3\baselineskip}\item Опровергни: если \(a^{2} = b^{2}\), то \(a = b\) для любых действительных \(a,\allowbreak b\).
\Needspace{3\baselineskip}\item Докажи по индукции: \(1 + 2 + \ldots  + n = \frac{n (n + 1)}{2}\) для натуральных \(n \ge  1\).
\Needspace{3\baselineskip}\item Докажи: если целое \(n\) нечётно, то \(n^{2}\) нечётно. Сформулируй равносильное утверждение через противоположное условие.
\Needspace{3\baselineskip}\item Из жизни. В кафе все комплексные обеды включают напиток. Покупатель получил напиток. Обязательно ли он заказал комплексный обед?
\Needspace{3\baselineskip}\item Из жизни. В группе \(20\) человек. Чай любят \(12\), кофе любят \(9\), оба напитка любят \(5\). Сколько любят хотя бы один из напитков и сколько не любят ни один?
\Needspace{3\baselineskip}\item Из жизни. Объявление гласит: «Все автобусы на этом маршруте низкопольные». Какое наблюдение опровергнет утверждение? Сформулируй его отрицание.
\Needspace{3\baselineskip}\item[\textbf{14.16.}\starmark] Из \(30\) учеников \(18\) изучают английский, \(16\) немецкий, каждый изучает хотя бы один из этих языков. Сколько изучают оба и сколько только английский?
\Needspace{3\baselineskip}\item[\textbf{14.17.}\starmark] Докажи: если квадрат целого числа делится на \(3\), то само число делится на \(3\).
\Needspace{3\baselineskip}\item[\textbf{14.18.}\starmark] Докажи по индукции: \(1 + 3 + 5 + \ldots  + (2 n - 1) = n^{2}\) для \(n \ge  1\).
\end{enumerate}
\clearpage\section*{Ответы и краткие решения}\addcontentsline{toc}{section}{Ответы и краткие решения}
\subsection*{1. Числа и свойства арифметических действий}
\begin{enumerate}[label=1.\arabic*.,leftmargin=2.8em,itemsep=5pt]
\item \(- 8 < - 3;\allowbreak 0 > - 5;\allowbreak - 4 < 4\).
\item \(- 12\).
\item \(3\).
\item \(- 3\).
\item \(- 8\).
\item \(7\). Делить на ноль нельзя.
\item \(- 14\).
\item \(- 20\).
\item \(- 26\).
\item \(- 4\). В записи \(- 2^{2}\) степень относится только к двойке.
\item \(- 2 {}^{\circ} C\).
\item На уровне \(2\).
\item \(115\) очков.
\item[\textbf{1.14.}\starmark] \(3\). Скобки равны \(8\) и \(3\), поэтому \(- 6 + 9 = 3\).
\item[\textbf{1.15.}\starmark] \((18 - 6) : (3 + 1) = 3\).
\item[\textbf{1.16.}\starmark] \(8\) верных и \(4\) неверных: \(8 \cdot  5 - 4 \cdot  3 = 28\).
\end{enumerate}
\subsection*{2. Делимость, простые числа, НОД и НОК}
\begin{enumerate}[label=2.\arabic*.,leftmargin=2.8em,itemsep=5pt]
\item \(1,\allowbreak 2,\allowbreak 3,\allowbreak 4,\allowbreak 6,\allowbreak 9,\allowbreak 12,\allowbreak 18,\allowbreak 36\).
\item \(1,\allowbreak 2,\allowbreak 3,\allowbreak 4,\allowbreak 6,\allowbreak 8,\allowbreak 12,\allowbreak 16,\allowbreak 24,\allowbreak 48\).
\item \(1,\allowbreak 2,\allowbreak 5,\allowbreak 10\).
\item \(21,\allowbreak 28,\allowbreak 35,\allowbreak 42,\allowbreak 49,\allowbreak 56\).
\item \(104\).
\item Простые: \(17,\allowbreak 31,\allowbreak 53\). Составные: \(25,\allowbreak 39,\allowbreak 49,\allowbreak 51\).
\item \(23,\allowbreak 29,\allowbreak 31,\allowbreak 37\).
\item Нет: такое натуральное число не меньше \(10\) и имеет делитель \(2\), отличный от \(1\) и самого числа.
\item Нет, например \(9 = 3 \cdot  3\).
\item На \(3\) и на \(9\) делятся \(126,\allowbreak 234,\allowbreak 405,\allowbreak 702;\allowbreak 517\) не делится ни на \(3\), ни на \(9\).
\item \(230,\allowbreak 480\).
\item \(2,\allowbreak 5,\allowbreak 8\).
\item \(0\) или \(9\).
\item \(0,\allowbreak 6\).
\item \(2^{3} \cdot  3^{2}\).
\item \(2 \cdot  3 \cdot  5^{2}\).
\item \(2^{2} \cdot  3^{2} \cdot  7\).
\item \(90 = 2 \cdot  3^{2} \cdot  5\).
\item НОД \(= 12\); НОК \(= 72\).
\item НОД \(= 9\); НОК \(= 90\).
\item НОД \(= 1\); НОК \(= 120\).
\item НОД \(= 6\); НОК \(= 180\).
\item \(12\) наборов, в каждом \(2\) карандаша и \(3\) ручки.
\item В 9:36.
\item Сторона \(12\) см; \(35\) квадратов.
\item \(1,\allowbreak 2,\allowbreak 3,\allowbreak 4,\allowbreak 6,\allowbreak 8,\allowbreak 12,\allowbreak 24\).
\item \(6,\allowbreak 12,\allowbreak 18,\allowbreak 24,\allowbreak 30\).
\item \(2,\allowbreak 13,\allowbreak 29\). Число \(1\) не простое и не составное.
\item На все перечисленные.
\item \(2^{2} \cdot  3 \cdot  7\).
\item НОД \(= 6\); НОК \(= 36\).
\item \(14\) подарков; по \(3\) яблока и \(4\) мандарина.
\item Через \(24\) секунды.
\item Нет: \(97\) — простое число.
\item[\textbf{2.35.}\starmark] \(73\): НОК\((6,8,\allowbreak 9) = 72\), затем прибавляем \(1\).
\item[\textbf{2.36.}\starmark] \(36\). Решение: \(\frac{6 \cdot 180}{30} = 36\).
\item[\textbf{2.37.}\starmark] \(3150\) и \(3555\). Сумма цифр равна соответственно \(9\) и \(18\).
\end{enumerate}
\subsection*{3. Обыкновенные дроби}
\begin{enumerate}[label=3.\arabic*.,leftmargin=2.8em,itemsep=5pt]
\item \(\frac{3}{4}\).
\item \(3 \frac{2}{5}\).
\item \(\frac{17}{7}\).
\item \(\frac{5}{8} > \frac{7}{12}\).
\item \(- \frac{1}{12}\).
\item \(\frac{31}{45}\).
\item \(- \frac{2}{3}\).
\item \(\frac{3}{2}\).
\item \(\frac{1}{2}\).
\item \(1\).
\item \(\frac{2}{3}\).
\item \(32\) страницы.
\item \(\frac{17}{12}\) кг, или \(1 \frac{5}{12}\) кг.
\item \(6\) стаканов.
\item \(9\) км.
\item[\textbf{3.16.}\starmark] \(-\frac{1}{4}\). Решение: \((-\frac{1}{12}) : (\frac{1}{9}) + \frac{1}{2} = -\frac{3}{4} + \frac{1}{2} = -\frac{1}{4}\).
\item[\textbf{3.17.}\starmark] \(180\) страниц: осталось \((\frac{2}{3}) \cdot  (\frac{3}{5}) = \frac{2}{5}\) книги.
\item[\textbf{3.18.}\starmark] \(30\) л: осталось \((\frac{3}{4}) \cdot  (\frac{2}{3}) = \frac{1}{2}\) бака.
\end{enumerate}
\subsection*{4. Десятичные дроби и проценты}
\begin{enumerate}[label=4.\arabic*.,leftmargin=2.8em,itemsep=5pt]
\item \(3,836\).
\item \(5,53\).
\item \(0,12\).
\item \(60\).
\item \(0,875\).
\item \(\frac{1}{8}\).
\item \(0{,}(6)\).
\item \(3,14;\allowbreak 3,142\).
\item \(36\).
\item \(180\).
\item На \(15 \%\).
\item \(960\) рублей; снижение на \(4 \%\).
\item \(225\) рублей.
\item \(4080\) рублей.
\item \(13,33\) км; после округления \(13,3\) км.
\item[\textbf{4.16.}\starmark] \(3000\) рублей: \(\frac{2160}{0,8 \cdot  0,9}\).
\item[\textbf{4.17.}\starmark] На \(20 \%\). Решение: \(\frac{25}{125} = 0{,}2 = 20 \%\).
\item[\textbf{4.18.}\starmark] \(16 \%\). Решение: \(\frac{20 + 60}{500} \cdot 100 \% = 16 \%\).
\end{enumerate}
\subsection*{5. Отношения и пропорции}
\begin{enumerate}[label=5.\arabic*.,leftmargin=2.8em,itemsep=5pt]
\item \(3 : 4\).
\item \(x = 12\).
\item \(x = 6\).
\item \(24\) и \(40\).
\item \(750\) г.
\item \(288\) рублей.
\item \(3,5\) км.
\item За \(4\) дня.
\item \(2\) часа.
\item \(0,6\) л сока; \(0,9\) л воды.
\item \(y = 12\); удвоится.
\item \(y = 4\); уменьшится вдвое.
\item \(750\) г.
\item За \(6\) часов.
\item \(0,25\) л сиропа и \(1,75\) л воды.
\item[\textbf{5.16.}\starmark] \(10\) дней: \(6 \cdot  8 \cdot  \frac{5}{4 \cdot  6}\).
\item[\textbf{5.17.}\starmark] \(7\) л: \(\frac{2 k}{5 k + 3} = \frac{1}{4}\) даёт \(k = 1\).
\item[\textbf{5.18.}\starmark] При равных временах \(75\) км/ч. При равных расстояниях \(72\) км/ч: \(\frac{2}{\frac{1}{60} + \frac{1}{90}}\).
\end{enumerate}
\subsection*{6. Степени с целым показателем}
\begin{enumerate}[label=6.\arabic*.,leftmargin=2.8em,itemsep=5pt]
\item \(32\).
\item \(81\) и \(- 81\).
\item \(\frac{9}{8}\).
\item \(5^{7}\).
\item \(2^{5}\).
\item \(3^{8}\).
\item \(\frac{9}{4}\).
\item \(100\).
\item \(5,6 \cdot  10^{- 4}\).
\item \(720000\).
\item \(a^{3}\).
\item \(- 1\).
\item \(2^{8} = 256\) бактерий.
\item \(10^{6}\) м\(^{2}\) \(= 1000000\) м\(^{2}\).
\item \(1,2 \cdot  10^{9}\) байт.
\item[\textbf{6.16.}\starmark] \(16 = 2^{4}\).
\item[\textbf{6.17.}\starmark] \(\frac{9}{2}\). Решение: \(\frac{3^{4} \cdot 2}{3^{2} \cdot 4} = \frac{9}{2}\).
\item[\textbf{6.18.}\starmark] Сначала \(2\) МБ; через \(6\) дней \(1458\) МБ.
\end{enumerate}
\subsection*{7. Корни и действительные числа}
\begin{enumerate}[label=7.\arabic*.,leftmargin=2.8em,itemsep=5pt]
\item \(9;\allowbreak - 3\).
\item \(\frac{7}{10}\).
\item \(6 \sqrt{2}\).
\item \(5 \sqrt{2}\).
\item \(\sqrt{45}\).
\item \(6\).
\item \(5\).
\item \(7\).
\item Между \(5\) и \(6\).
\item \(\sqrt{7} < 2,7\), так как \(7 < 2,7^{2} = 7,29\).
\item \(\sqrt{16},\allowbreak - \frac{3}{7},\allowbreak 0{,}(3)\).
\item \(\sqrt{- 9}\) не существует; \(\sqrt[3]{- 8} = - 2\).
\item \(7\) м.
\item \(5\) м.
\item \(0,5\) м, или \(50\) см.
\item[\textbf{7.16.}\starmark] \(2\sqrt{3}\). Решение: \(4\sqrt{3} + 5\sqrt{3} - 7\sqrt{3} = 2\sqrt{3}\).
\item[\textbf{7.17.}\starmark] Квадрат левой части \(5 + 2 \sqrt{6} < 10\), поскольку \(\sqrt{6} < \frac{5}{2}\). Поэтому исходное утверждение неверно: \(\sqrt{2} + \sqrt{3} < \sqrt{10}\).
\item[\textbf{7.18.}\starmark] \(4\) м и \(12\) м; диагональ \(4 \sqrt{10}\) м.
\end{enumerate}
\subsection*{8. Модуль и числовые промежутки}
\begin{enumerate}[label=8.\arabic*.,leftmargin=2.8em,itemsep=5pt]
\item \(11\).
\item \(3\).
\item \(11\).
\item \(| - 6 | > | 5 |\).
\item \((- 2;\allowbreak 3]\).
\item \(- 4 \le  x < 1\).
\item \(- 2,\allowbreak - 1,\allowbreak 0,\allowbreak 1,\allowbreak 2\).
\item \((1;\allowbreak 4]\).
\item \((- \infty ;\allowbreak 5)\).
\item \(x = - 5\) или \(x = 5\).
\item \(- 1 < x < 5\).
\item \(| x |;\allowbreak 4\).
\item На \(40 {}^{\circ} C\).
\item \(| x - 50 | \le  0,2;\allowbreak [49,8;\allowbreak 50,2]\) мм.
\item \(120 \le  h < 190;\allowbreak [120;\allowbreak 190)\); при \(190\) см нельзя.
\item[\textbf{8.16.}\starmark] \(x = - 2\) или \(x = 4\). Между \(- 1\) и \(3\) сумма равна \(4\).
\item[\textbf{8.17.}\starmark] \(- 2,\allowbreak - 1,\allowbreak 2,\allowbreak 3\): допустимый отрезок \([- 2;\allowbreak 3]\), исключаются \(0\) и \(1\).
\item[\textbf{8.18.}\starmark] \((20;\allowbreak 23]\): пересечение \([17;\allowbreak 23]\) и \((20;\allowbreak 28)\).
\end{enumerate}
\subsection*{9. Преобразование алгебраических выражений}
\begin{enumerate}[label=9.\arabic*.,leftmargin=2.8em,itemsep=5pt]
\item \(4 x - 17\).
\item \(6 a + 2 b\).
\item \(x^{2} + 6 x + 9\).
\item \(4 x^{2} - 1\).
\item \(3 x (2 x - 3)\).
\item \((x - 4) (x + 4)\).
\item \((x + 2) (x + 3)\).
\item \(x + 1\) при \(x \ne  1\).
\item \(\frac{7}{2 x},\allowbreak x \ne  0\).
\item \(\frac{a}{b},\allowbreak a \ne  0\) и \(b \ne  0\).
\item \(\frac{x - 2}{x + 2},\allowbreak x \ne  - 2\).
\item \(\sqrt{5} + 2\).
\item \(3 x + 2 (x - 100) = 5 x - 200\) рублей.
\item \(x^{2} + 7 x + 10\) м\(^{2}\).
\item \((x + 3)^{2} - x^{2} = 6 x + 9\) м\(^{2}\).
\item[\textbf{9.16.}\starmark] \(\frac{x + 3}{x + 2},\allowbreak x \ne  3\) и \(x \ne  - 2\).
\item[\textbf{9.17.}\starmark] \(1\) при \(x \ne  - 1\) и \(x \ne  1\).
\item[\textbf{9.18.}\starmark] Разность равна \(4 ab\); при данных значениях \(39964\).
\end{enumerate}
\subsection*{10. Уравнения, системы и неравенства}
\begin{enumerate}[label=10.\arabic*.,leftmargin=2.8em,itemsep=5pt]
\item \(x = 6\).
\item \(x = - 5\).
\item \(x = 2\) или \(3\).
\item \(x = - 2\) или \(2\).
\item \(x = 5\); ОДЗ: \(x \ne  2\).
\item \(x = 4,\allowbreak y = 3\).
\item \(x < - 4\).
\item \([- 3;\allowbreak 3]\).
\item \((- \infty ;\allowbreak - 2) \cup  (1;\allowbreak + \infty )\).
\item \((- 1;\allowbreak 3]\).
\item \(x = - 2\) или \(3\).
\item \(x=2\). ОДЗ: \(x\ge-2\). Для выполнения равенства необходимо \(x\ge0\). После возведения в квадрат кандидаты \(2\) и \(-1\); проверка исключает \(-1\).
\item \(12\) км.
\item \(120\) рублей.
\item \(5\) билетов.
\item[\textbf{10.16.}\starmark] \(x = 3\). После возведения в квадрат корни \(3\) и \(- 1;\allowbreak - 1\) посторонний.
\item[\textbf{10.17.}\starmark] \(x < - 3\): дробь после переноса равна \(- \frac{5}{x + 3}\).
\item[\textbf{10.18.}\starmark] \(5\) м и \(12\) м: \(x (17 - x) = 60\).
\end{enumerate}
\subsection*{11. Функции и графики}
\begin{enumerate}[label=11.\arabic*.,leftmargin=2.8em,itemsep=5pt]
\item \(- 7,\allowbreak - 3,\allowbreak 3\).
\item Все действительные \(x\), кроме \(4\).
\item \([- 3;\allowbreak + \infty )\).
\item \(x = - 2\) и \(2\).
\item При \(x > 2\).
\item С осью \(y\): \((0;\allowbreak 1)\); с осью \(x\): \((-\frac{1}{2};\allowbreak 0)\).
\item Вершина \((2;\allowbreak - 1)\); сдвиг \(y = x^{2}\) на \(2\) вправо и на \(1\) вниз.
\item Убывает на \((- \infty ;\allowbreak 0]\), возрастает на \([0;\allowbreak + \infty )\).
\item \(f\) нечётная; \(g\) чётная.
\item \((x + 1)^{2};\allowbreak x^{2} + 1\).
\item \([2;\allowbreak + \infty )\).
\item \((- 1;\allowbreak 1)\) и \((2;\allowbreak 4)\).
\item \(C (0) = 200,\allowbreak C (2) = 500,\allowbreak C (4) = 800\) рублей; луч из точки \((0;\allowbreak 200)\).
\item \(V (t) = 60 - 5 t;\allowbreak 0 \le  t \le  12\); бак пуст через \(12\) минут.
\item Равны при \(5\) км; при \(2\) км дешевле \(B\): \(80\) против \(140\); при \(8\) км дешевле \(A\): \(260\) против \(320\) рублей.
\item[\textbf{11.16.}\starmark] \([1;\allowbreak 3) \cup  (3;\allowbreak + \infty )\).
\item[\textbf{11.17.}\starmark] Вершина \((2;\allowbreak 1);\allowbreak x \in  [- 1;\allowbreak 5]\).
\item[\textbf{11.18.}\starmark] \(S = x (10 - x) = 25 - (x - 5)^{2},\allowbreak 0 < x < 10\). Максимум \(25\) м\(^{2}\) при сторонах \(5\) м.
\end{enumerate}
\subsection*{12. Степенные, показательные и логарифмические функции}
\begin{enumerate}[label=12.\arabic*.,leftmargin=2.8em,itemsep=5pt]
\item Обе: все действительные числа, кроме \(0\); ветви в \(I\) и \(III\) четвертях.
\item Обе: \([0;\allowbreak + \infty )\); график проходит через \((0;\allowbreak 0),\allowbreak (1;\allowbreak 1),\allowbreak (4;\allowbreak 2)\).
\item \(x = - 3\).
\item \(x = 5\).
\item \(x = 3\).
\item \(x < 3\).
\item \(3\).
\item \(- 2\).
\item \(x = 9\); ОДЗ: \(x > 1\).
\item \(x < 5\).
\item \(3\).
\item Первое меньше второго в первой паре; первое больше второго во второй паре.
\item \(10000\) рублей.
\item Через \(4\) часа.
\item \(10\) шагов в заданной модели.
\item[\textbf{12.16.}\starmark] \(x = 0\) или \(x = 2\). Подстановка \(t = 2^{x}\) даёт \(t^{2} - 5t + 4 = 0\), откуда \(t = 1\) или \(t = 4\).
\item[\textbf{12.17.}\starmark] \(x = 3\): ОДЗ \(x > 1;\allowbreak x^{2} - 1 = 8\); отрицательный корень исключается.
\item[\textbf{12.18.}\starmark] \(3\) года: \(1,1^{3} = 1,331\), а через \(2\) года сумма \(12100\) рублей.
\end{enumerate}
\subsection*{13. Тригонометрия}
\begin{enumerate}[label=13.\arabic*.,leftmargin=2.8em,itemsep=5pt]
\item \(\frac{\pi }{6},\allowbreak \frac{\pi }{2},\allowbreak \frac{5 \pi }{6}\).
\item \(45 {}^{\circ}\) и \(210 {}^{\circ}\).
\item \(\frac{1}{2};\allowbreak \frac{1}{2};\allowbreak 1\).
\item \(\sin \) положителен, \(\cos \) отрицателен.
\item \(- \frac{1}{2};\allowbreak \frac{1}{2}\).
\item \(3\).
\item \(\cos  x = \frac{4}{5};\allowbreak \tan  x = \frac{3}{4}\).
\item \(\sin  x;\allowbreak - \cos  x\).
\item \(0,\allowbreak \pi ,\allowbreak 2 \pi \).
\item \(\frac{\pi }{3},\allowbreak \frac{5 \pi }{3}\).
\item Нули: \(0,\allowbreak \pi ,\allowbreak 2 \pi \); максимум \(1\) при \(\frac{\pi }{2}\); минимум \(- 1\) при \(\frac{3 \pi }{2}\).
\item \(\pi \).
\item \(6\) м.
\item \(5\) м.
\item \(12\) м, \(22\) м, \(2\) м соответственно.
\item[\textbf{13.16.}\starmark] \(\frac{\pi}{6},\allowbreak \frac{\pi}{2},\allowbreak \frac{5\pi}{6}\). После решения квадратного уравнения получаем \(\sin x = \frac{1}{2}\) или \(\sin x = 1\).
\item[\textbf{13.17.}\starmark] \(4 \sqrt{3}\) м.
\item[\textbf{13.18.}\starmark] Минимум \(3\) м, максимум \(27\) м; первый максимум при \(t = 20\) с.
\end{enumerate}
\subsection*{14. Множества, логика и доказательства}
\begin{enumerate}[label=14.\arabic*.,leftmargin=2.8em,itemsep=5pt]
\item Пересечение: \(\{3\}\); объединение: \(\{1,\allowbreak 2,\allowbreak 3,\allowbreak 4\}\).
\item \(\forall  x \in  \mathbb{R} : x^{2} \ge  0\).
\item Существует целое \(x\), квадрат которого равен \(9\). Верно: \(x = 3\) или \(- 3\).
\item Нет: при \(x = 0\) квадрат равен \(0\).
\item В наборе есть хотя бы одно число, которое не положительно \((\le  0)\).
\item Прямая верна; обратная нет, контрпример \(6\).
\item Необходимое, но недостаточное: \(2\) чётно, но на \(6\) не делится.
\item \(2 m + 2 n = 2 (m + n)\), где \(m,\allowbreak n \in  \mathbb{Z}\).
\item Пусть \(n=2k+1\), где \(k\in\mathbb Z\). Тогда \(n^2=2(2k^2+2k)+1\); число \(2k^2+2k\) целое, поэтому \(n^2\) нечётно.
\item Контрпример: \(a = 1,\allowbreak b = - 1\).
\item База \(n = 1\) верна. Если сумма до \(k\) равна \(\frac{k (k + 1)}{2}\), то после прибавления \(k + 1\) получаем \((k + 1) \frac{k + 2}{2}\).
\item Доказательство как в задаче \(9\). Контрапозиция: если \(n^{2}\) чётно, то \(n\) чётно.
\item Нет: напиток мог быть куплен отдельно.
\item Хотя бы один любят \(16\); ни один не любят \(4\).
\item Достаточно одного автобуса этого маршрута, который не низкопольный. Отрицание: существует такой автобус.
\item[\textbf{14.16.}\starmark] Оба языка изучают \(4\); только английский \(14\).
\item[\textbf{14.17.}\starmark] При делении любого целого числа на \(3\) возможны остатки \(0,1,2\). Остатки квадратов: \(0,\allowbreak 1,\allowbreak 1\). Нулевой остаток возможен только при остатке \(0\) у самого числа.
\item[\textbf{14.18.}\starmark] База: \(1 = 1^{2}\). Переход: \(k^{2} + (2 k + 1) = (k + 1)^{2}\).
\end{enumerate}
\end{document}
